% Options for packages loaded elsewhere
\PassOptionsToPackage{unicode}{hyperref}
\PassOptionsToPackage{hyphens}{url}
\documentclass[
]{article}
\usepackage{xcolor}
\usepackage{amsmath,amssymb}
\setcounter{secnumdepth}{-\maxdimen} % remove section numbering
\usepackage{iftex}
\ifPDFTeX
  \usepackage[T1]{fontenc}
  \usepackage[utf8]{inputenc}
  \usepackage{textcomp} % provide euro and other symbols
\else % if luatex or xetex
  \usepackage{unicode-math} % this also loads fontspec
  \defaultfontfeatures{Scale=MatchLowercase}
  \defaultfontfeatures[\rmfamily]{Ligatures=TeX,Scale=1}
\fi
\usepackage{lmodern}
\ifPDFTeX\else
  % xetex/luatex font selection
\fi
% Use upquote if available, for straight quotes in verbatim environments
\IfFileExists{upquote.sty}{\usepackage{upquote}}{}
\IfFileExists{microtype.sty}{% use microtype if available
  \usepackage[]{microtype}
  \UseMicrotypeSet[protrusion]{basicmath} % disable protrusion for tt fonts
}{}
\makeatletter
\@ifundefined{KOMAClassName}{% if non-KOMA class
  \IfFileExists{parskip.sty}{%
    \usepackage{parskip}
  }{% else
    \setlength{\parindent}{0pt}
    \setlength{\parskip}{6pt plus 2pt minus 1pt}}
}{% if KOMA class
  \KOMAoptions{parskip=half}}
\makeatother
\usepackage{color}
\usepackage{fancyvrb}
\newcommand{\VerbBar}{|}
\newcommand{\VERB}{\Verb[commandchars=\\\{\}]}
\DefineVerbatimEnvironment{Highlighting}{Verbatim}{commandchars=\\\{\}}
% Add ',fontsize=\small' for more characters per line
\newenvironment{Shaded}{}{}
\newcommand{\AlertTok}[1]{\textcolor[rgb]{1.00,0.00,0.00}{\textbf{#1}}}
\newcommand{\AnnotationTok}[1]{\textcolor[rgb]{0.38,0.63,0.69}{\textbf{\textit{#1}}}}
\newcommand{\AttributeTok}[1]{\textcolor[rgb]{0.49,0.56,0.16}{#1}}
\newcommand{\BaseNTok}[1]{\textcolor[rgb]{0.25,0.63,0.44}{#1}}
\newcommand{\BuiltInTok}[1]{\textcolor[rgb]{0.00,0.50,0.00}{#1}}
\newcommand{\CharTok}[1]{\textcolor[rgb]{0.25,0.44,0.63}{#1}}
\newcommand{\CommentTok}[1]{\textcolor[rgb]{0.38,0.63,0.69}{\textit{#1}}}
\newcommand{\CommentVarTok}[1]{\textcolor[rgb]{0.38,0.63,0.69}{\textbf{\textit{#1}}}}
\newcommand{\ConstantTok}[1]{\textcolor[rgb]{0.53,0.00,0.00}{#1}}
\newcommand{\ControlFlowTok}[1]{\textcolor[rgb]{0.00,0.44,0.13}{\textbf{#1}}}
\newcommand{\DataTypeTok}[1]{\textcolor[rgb]{0.56,0.13,0.00}{#1}}
\newcommand{\DecValTok}[1]{\textcolor[rgb]{0.25,0.63,0.44}{#1}}
\newcommand{\DocumentationTok}[1]{\textcolor[rgb]{0.73,0.13,0.13}{\textit{#1}}}
\newcommand{\ErrorTok}[1]{\textcolor[rgb]{1.00,0.00,0.00}{\textbf{#1}}}
\newcommand{\ExtensionTok}[1]{#1}
\newcommand{\FloatTok}[1]{\textcolor[rgb]{0.25,0.63,0.44}{#1}}
\newcommand{\FunctionTok}[1]{\textcolor[rgb]{0.02,0.16,0.49}{#1}}
\newcommand{\ImportTok}[1]{\textcolor[rgb]{0.00,0.50,0.00}{\textbf{#1}}}
\newcommand{\InformationTok}[1]{\textcolor[rgb]{0.38,0.63,0.69}{\textbf{\textit{#1}}}}
\newcommand{\KeywordTok}[1]{\textcolor[rgb]{0.00,0.44,0.13}{\textbf{#1}}}
\newcommand{\NormalTok}[1]{#1}
\newcommand{\OperatorTok}[1]{\textcolor[rgb]{0.40,0.40,0.40}{#1}}
\newcommand{\OtherTok}[1]{\textcolor[rgb]{0.00,0.44,0.13}{#1}}
\newcommand{\PreprocessorTok}[1]{\textcolor[rgb]{0.74,0.48,0.00}{#1}}
\newcommand{\RegionMarkerTok}[1]{#1}
\newcommand{\SpecialCharTok}[1]{\textcolor[rgb]{0.25,0.44,0.63}{#1}}
\newcommand{\SpecialStringTok}[1]{\textcolor[rgb]{0.73,0.40,0.53}{#1}}
\newcommand{\StringTok}[1]{\textcolor[rgb]{0.25,0.44,0.63}{#1}}
\newcommand{\VariableTok}[1]{\textcolor[rgb]{0.10,0.09,0.49}{#1}}
\newcommand{\VerbatimStringTok}[1]{\textcolor[rgb]{0.25,0.44,0.63}{#1}}
\newcommand{\WarningTok}[1]{\textcolor[rgb]{0.38,0.63,0.69}{\textbf{\textit{#1}}}}
\usepackage{longtable,booktabs,array}
\newcounter{none} % for unnumbered tables
\usepackage{calc} % for calculating minipage widths
% Correct order of tables after \paragraph or \subparagraph
\usepackage{etoolbox}
\makeatletter
\patchcmd\longtable{\par}{\if@noskipsec\mbox{}\fi\par}{}{}
\makeatother
% Allow footnotes in longtable head/foot
\IfFileExists{footnotehyper.sty}{\usepackage{footnotehyper}}{\usepackage{footnote}}
\makesavenoteenv{longtable}
\setlength{\emergencystretch}{3em} % prevent overfull lines
\providecommand{\tightlist}{%
  \setlength{\itemsep}{0pt}\setlength{\parskip}{0pt}}
% Header includes for the Puppeteer Papers pandoc + tectonic pipeline.
% Every line here is load-bearing; the reason is stated so it is not dropped later.

% pandoc emits \xmpquote{} inside \hypersetup{pdfkeywords} when the YAML has
% keywords:. Without hyperxmp tectonic dies "Undefined control sequence \xmpquote".
\usepackage{hyperxmp}

% Long monospace paths inside Appendix A's tables overflow the column. seqsplit
% breaks them. SCOPED TO LONGTABLES ON PURPOSE: a global \renewcommand\texttt with
% seqsplit dies on inline code containing escaped specials (\_ etc.).
\usepackage{seqsplit}
\usepackage{etoolbox}
\AtBeginEnvironment{longtable}{\let\origtexttt\texttt\renewcommand\texttt[1]{\origtexttt{\seqsplit{#1}}}}
\AtEndEnvironment{longtable}{\let\texttt\origtexttt}

% Wrap long lines inside highlighted (fenced) code blocks instead of overflowing.
\usepackage{fvextra}
\DefineVerbatimEnvironment{Highlighting}{Verbatim}{breaklines,breakanywhere,commandchars=\\\{\}}

% Let a long inline \texttt move to the next line rather than overflow the margin.
\AtBeginDocument{\sloppy}

% tectonic's text font silently DROPS these to U+FFFD ("~1.3 ms" losing its ~).
% newunicodechar catches text-mode characters only, not math-mode ones.
\usepackage{newunicodechar}
\newunicodechar{≈}{\ensuremath{\approx}}
\newunicodechar{∀}{\ensuremath{\forall}}
\newunicodechar{∈}{\ensuremath{\in}}
\newunicodechar{↔}{\ensuremath{\leftrightarrow}}
\newunicodechar{α}{\ensuremath{\alpha}}
\newunicodechar{β}{\ensuremath{\beta}}
\newunicodechar{▲}{\ensuremath{\blacktriangle}}
\newunicodechar{▼}{\ensuremath{\blacktriangledown}}
\newunicodechar{►}{\ensuremath{\blacktriangleright}}
\newunicodechar{◄}{\ensuremath{\blacktriangleleft}}
\newunicodechar{₁}{\textsubscript{1}}
\newunicodechar{₂}{\textsubscript{2}}
\newunicodechar{ₙ}{\textsubscript{n}}
\newunicodechar{ᵢ}{\textsubscript{i}}
% U+21B3 is the marker Paper 9's Table 3 uses for "consequence of the row above",
% and its caption refers to the symbol by name — so dropping it empties three cells
% AND orphans the caption. Neither Latin Modern nor Consolas has it; Cambria Math
% does, and taking the glyph from there keeps the author's mark rather than
% substituting \hookrightarrow for it.
\newfontfamily\arrowfont{Cambria Math}
\newunicodechar{↳}{{\arrowfont ↳}}

% The default mono (Latin Modern) lacks U+2500 box drawing, so ASCII diagrams in
% code blocks drop out. Consolas has them. System-specific, hence the build's
% "may not be reproducible" warning — acceptable for this Windows pipeline.
\setmonofont{Consolas}[Scale=MatchLowercase]
\usepackage{bookmark}
\IfFileExists{xurl.sty}{\usepackage{xurl}}{} % add URL line breaks if available
\urlstyle{same}
\hypersetup{
  pdftitle={Identity Precedes Staging: one play{,} many stages},
  pdfauthor={Alvaro Rivera},
  pdfkeywords={\xmpquote{identity}, \xmpquote{staging}, \xmpquote{domain
independence}, \xmpquote{decoupling}, \xmpquote{ports and
adapters}, \xmpquote{projection}, \xmpquote{narrative
recognition}, \xmpquote{journaled systems}, \xmpquote{actor-native
architecture}, \xmpquote{design theory}, \xmpquote{puppeteer
framework}},
  hidelinks,
  pdfcreator={LaTeX via pandoc}}

\title{Identity Precedes Staging: one play, many stages}
\author{Alvaro Rivera}
\date{2026-08-07}

\begin{document}
\maketitle
\begin{abstract}
A program's domain is usually written together with the code that
deploys it, so a new place to run, or a new kind of client, means
editing the domain. This paper reports one that was not. A game's rules
--- holding no input or output, naming no framework in its build graph
--- ran on five deployment configurations and were read by six clients,
the last of them across three machines, at zero edits to the domain and
zero test doubles in its own tests. A deployment is called a
\emph{staging} throughout, after theatre: one script, many productions.

The mechanisms that allow this are old and none is claimed. What is
claimed is a property of the domain, measured against an orthodox
ports-and-adapters version of the same game built for the comparison.
That baseline matches the zero per staging. It \emph{decouples} a domain
where this arrangement \emph{closes} one, on two independent
measurements: the domain declares no obligation outward, where the
baseline declares three driven ports --- whence its stand-ins, its
public surface, and the state it admits from outside --- and its
reconstitution needs no surface of its own, where the baseline grew one
at 56 lines added and 5 removed. Calling those two \emph{closure}, and
reading the invariance as an \emph{identity} prior to any staging, are
readings and not further measurements. The paper is analytic in Gregor's
(2006, Type I) sense, and offers a question to put to a system already
built rather than a prescription. Two limits bound it: the domain is a
pure computational core that causes no external effect, and its author
commissioned the work.
\end{abstract}

\section{Identity Precedes Staging: one play, many
stages}\label{identity-precedes-staging-one-play-many-stages}

\subsection{TL;DR}\label{tldr}

A domain and its deployment are usually written together, so a new place
to run means editing the domain. They can be separated completely, and
this paper measures one that is. A Tetris board was held fixed while
five stagings and six clients changed around it: zero domain edits, zero
test doubles. The domain declares no interface for what drives it or for
what it emits, so the contract that would be a port is provided by the
substrate and bound in the staging. None of that geometry is new. What
the paper adds is a measurement of it, against a ports-and-adapters
version of the same game built for the comparison --- which matches the
zero per staging. What it does not match is of another kind: ports and
adapters \emph{decouples} a domain, while this arrangement \emph{closes}
one. Two measurements: a ported domain declares obligations outward ---
three driven ports, and with them three stand-ins without which twenty
of sixty-four tests do not run, an operation per port on its public
surface, and a state it will accept from outside --- and it needs a
reconstitution surface of its own, costing 56 lines. This one declares
none and needs none, replay re-entering by the operations the acts were
performed through. Calling those two \emph{closure} is a reading of them
rather than a third count. Two findings arrived unsought, neither of
them this paper's result and both bounding how it may be read: an
incomplete record answers plausibly rather than failing loudly, so carry
the recorded state and assert a replay against it; and a fact a domain
emits must be derivable inside one actor, which bounds how a domain may
be cut.

\emph{Dependencies. This paper is part of the Puppeteer Papers, a series
of self-deposited preprints, and rests on four of them: the actor's
speech and \texttt{tell} (Paper 4), \texttt{Reaction} read as the
recognition of a routine (Paper 3), the server read as an accidental
category, ephemeral after bootstrap (Paper 7), and testimony (Paper 8),
whose observer receives an account it is told. Paper 8 noted, without
pursuing it, that a narration received is not yet a narrative
recognized; this paper takes up the recognition, and reaches it not as
its subject but as a consequence of the identity it argues.
Methodologically it is an analytic theory contribution in the sense of
Gregor's (2006) theory for analyzing (Type I): identity and staging are
constructs by which an architecture may be described and compared, and
the paper offers no prescription for building one --- its instrument is
a question that can be asked of a system already built (§10), not a rule
for making one. The labs of Appendix A are accordingly an existence
proof of realizability, not a design-science evaluation of cost or
benefit, and where the text has slipped into evaluative language the
slip is corrected rather than defended. Stated in the vocabulary of
software-engineering research method rather than of information systems:
the contribution is a descriptive model whose validation is by example
--- a system built and measured --- in Shaw's (2003) classification of
contribution and validation types. And in the terms of Stol and
Fitzgerald's (2018) framework, it trades generalizability over actors
and realism of context away entirely, being one domain on one framework,
in exchange for a precise and reproducible measurement of a single
quantity.}

\subsection{1. The Domain and the
Deployment}\label{the-domain-and-the-deployment}

In most systems, the code that defines what a system does is not cleanly
separable from the code that defines where it runs and which clients it
serves. Moving a domain from a console to a web server means editing it;
adding a second kind of client means adding to the thing the client
observes; splitting one process into several is treated,
uncontroversially, as building a different system. Where a system runs,
and whom it serves, are treated as properties of the system itself ---
so that a monolith and its distributed successor are counted as two
systems rather than one system deployed two ways.

There is a field in which this separation is routine rather than
aspirational: theatre. A play is written once and staged in many venues,
for many audiences, without being rewritten, and no one counts the
Broadway production and the school-gymnasium production as different
plays. The term is taken from there: throughout this paper a
\emph{staging} is a deployment --- a particular place a domain runs and
a particular client it serves --- and the domain is what is staged. Two
units are counted in what follows and they are not the same: a
\textbf{staging} is a deployment configuration, of which the example has
five, and a \textbf{host} is an executable project, of which it has
twelve, eleven reaching the domain through the actor and one writing
bytes to a pipe and reaching nothing. Counts of hosts are counts of
projects that had to compile and run; counts of stagings are counts of
arrangements the domain was placed in. It is a naming convenience, not
an argument, and earns its keep only if the separation it names can be
made real and measured --- which is the rest of the paper.

The claim is not that the domain is \emph{decoupled} from its
deployment, which is ordinary advice. It is that the domain is
\textbf{closed}: it declares no obligation for a deployment to meet, so
the dependency between the two runs one way only, and that direction is
a checkable property of the built system. The domain compiles and runs
with no staging bound to it. Every staging refers to the domain; the
domain refers to no staging. In that precise sense the domain comes
first: a staging is constructed against a domain that already exists and
does not know of it. This is all the title means by \emph{precedes} ---
a direction of dependence in the build, not a claim about time or
metaphysics --- and §2 measures it directly.

This is not dependency inversion under another name, though the
difference takes some care to state and §9 states it. Ports-and-adapters
architecture keeps adapters outside the domain, and for whatever a
domain needs \emph{from} the world it has the domain declare an
interface and depend on it. In the model examined here the domain
declares no interface at all --- not for what drives it, nor for what it
emits. It is entered by calls on its own operations, and it emits facts
under logical names, while where those facts go, to whom, and what
drives them are bound entirely outside it. Whether that is
\emph{cheaper} than the ported alternative is a separate question, and
mostly it is not: a ported version of this game was built and counted,
and a new transport costs its domain nothing either (§9, Appendix A).
What the ported design still has, and this one does not, is an open
obligation --- an interface it declared and depends on someone to
supply. That difference, not a cost, is what §9 measures.

Placed in the series, this is a question of a specific kind. Paper 4
asked who may speak; Paper 8 asked who may decide what an output is; the
series puts its questions in that form. This paper asks a prior one ---
not who is entitled, but what an entitlement is attributed to, since a
subject that does not hold still across its stagings cannot bear one at
all. It is a question about the identity of a domain, not about
infrastructure or topology. A later paper in the series will need that
question answered before it can treat a set of domains as a reusable
repertoire; that is not this paper's concern, which is the narrower,
checkable claim that a domain's identity is prior to and independent of
its staging. One consequence is developed in §6 and §7: when a domain's
acts are recorded as they are performed, the sequence of what it did is
available whole from its journal --- read from that record, wherever a
copy of it is kept, rather than assembled out of fragments held
separately.

One question follows, and a second that the first makes unavoidable. The
first (§2, §5): does anything actually hold a domain fixed while its
staging changes --- across both where it runs and which client it serves
--- can that be measured rather than asserted, and what does the
invariance reveal? The second (§3--§4, §6) is not a parallel
contribution but what the first needs in order to be legible at all:
since a client never reads a domain directly, through what does it read,
and does it read the domain's present or its past? An invariance nobody
can observe is not a result, so the reading has to be accounted for;
that is the whole of the second question's standing here, and §6 says
where it stops.

One thing about the order of the work belongs here too, because it
decides how the rest is to be read. \textbf{The arrangement preceded its
account; the experiments were built to expose and bound it.} The domain,
the actor and most of the stagings of Experiment A existed before
\emph{identity}, \emph{staging} and \emph{closure} were words for
anything here, and none of them was built to satisfy a construct. The
labs are the other half: the ported baseline, the three containers, and
the growth of §8.2 were made for this argument, and §8.3 says plainly
where that weakens what they show. So the constructs are instruments for
reading a construction rather than a design derived from a theory, and
the question they are answerable to is whether they describe faithfully
something already built --- which is the sense in which the labs are
validation by example and not a test of a hypothesis.

\subsection{2. Two Experiments}\label{two-experiments}

The claim of §1 is checkable, and this section checks it against a
running example: a Tetris game, whose board and rules are a small but
complete domain, written as a \texttt{Well} and a set of operations over
it. The domain is held fixed, and two things are varied around it in
turn --- where it runs, and which client it serves. The measurement in
each case is the same: the number of changes the variation forces on the
domain source. The evidence is the git history of the example
repository, and the anchors are collected in Appendix A.

The domain's independence has a concrete form before anything runs ---
the shape of the build. Every executable host in the example reaches the
domain through a single actor, and the domain reaches back to nothing
but the language's own library:

\begin{verbatim}
  11 hosts           -> TetrisActor
     TetrisAi, TetrisConsole, TetrisObserver, TetrisServer, TetrisStage,
     TetrisStageDuo, TetrisStageDuoTls, TetrisStageServer, TetrisWatch,
     TetrisWeb, TetrisWebRest

  TetrisActor        -> TetrisDomain, Puppeteer, Choreography
  TetrisDomain.Tests -> TetrisDomain

  TetrisDomain       -> (none)
\end{verbatim}

\textbf{Figure 1.} Every declared dependency between the example's
projects, read out of the project files on its main branch rather than
drawn, and grouped so that the direction is visible: an arrow means
\emph{this project is built against that one}. Every staging depends on
the same domain; the domain depends on none. Fourteen of the fifteen
projects appear --- \texttt{TetrisSend}, the pipe writer, is omitted
because it references nothing and is referenced by nothing, being a
separate executable that only writes bytes. Read downward, every arrow
leads toward \texttt{TetrisDomain} and none leads away from it.

The figure answers one question, \emph{who depends on whom}, and it is
worth saying what it does not answer. It does not show that the domain
leaves nothing open --- no port declared, no contract, no persistence or
transport requested. That is a different property, it is the stronger of
the two, and it is measured separately against a built alternative in
§9. A reader unfamiliar with this language's build files needs only the
arrow: a project that names another cannot be compiled without it, and a
project that names none can be compiled alone.

Declaring no reference is not the same as naming nothing, and one
qualification belongs with the figure rather than after it. Three hosts
do name a domain type, reaching it through the actor's reference, which
in this language is transitive: each writes
\texttt{typeof(TetrisDomain).Assembly} to hand the assembly to the
framework (\texttt{sm-duo-tls/Program.cs:33}, and the same line in
\texttt{sm-server} and \texttt{sm-duo}). What they name is an empty
public anchor and nothing else, so a staging holds a name and not a
handle. Lab E reads that fence at the level of members and reports its
two authored exceptions.

What the domain does instead is \emph{satisfy conventions}, which is a
different relation from declaring a dependency and the difference is
load-bearing. The framework finds the domain by reflection, so its
operations must be ordinary methods with coercible parameters, and the
empty public anchor is there so a host can name the assembly in its own
source --- a convenience, not a requirement, as Lab F reports. What
separates conformance from dependency is not interpretive: the domain
compiles and its own tests run with the framework absent from the build
graph entirely (Lab E). A dependency cannot be absent from the build and
still be met; a convention can.

That direction is what §1 called \emph{precede}, and one note on
vocabulary avoids confusion later. With arrows running from a dependent
to what it depends on, the domain is a \emph{sink}: every staging's
references lead to it and none leads out. Reverse the arrows and it is
the \emph{root}, every staging reachable from it and none from another.
Where later sections call the domain the stagings' common ancestor, that
reversed orientation is what is meant; §5 says what the property does
and does not establish.

\textbf{Experiment A --- change the stage.} The same \texttt{Well} is
run on a console; in a browser, with input and output carried over a
real WebSocket connection, and in a second browser variant over HTTP
with server-sent events; across two co-hosted actors coordinated by a
StageManager, once in memory and once over a real TLS channel; and
across three separate Docker containers, each a peer joined to the
others over real container-to-container TLS. These are genuinely
different deployments --- different processes, machines, transports, and
wire formats. The domain source is identical across all of them: a diff
of the domain directory between the first staging and the last comes
back empty (Appendix A). What changes from one staging to the next is
host code --- the shell that binds an input source and an output sink to
the actor. The example's own name for that shell is exact: the host is
``an accidental shell'' (\texttt{actor/IGameHost.cs:14}), and the same
\texttt{Well} runs on \texttt{PerformanceV2} or on \texttt{StageV2}
without knowing which. The example builds clean and its domain tests
pass, unchanged, in every staging.

A distributed demonstration invites more credit than it earns, so its
boundaries belong in plain sight. Three containers run the same
\texttt{Well} as three StageManager peers over genuine
container-to-container TLS; a scripted driver on the Director plays the
game, and the two casts converge to a byte-identical board. Adding that
staging left the diff of both the domain and the actor directories
empty. What is \emph{not} claimed is any test of resilience: no peer was
killed and no partition induced, so partial failure is untouched (§9).
The setup, the bootstrap, and what each node is handed are in Lab B.

\textbf{Experiment B --- change the client.} With the stage held at one
host, the client is varied instead. The same board is played by a human
at a keyboard reading an on-screen grid; by an automated player that
sends moves over a pipe and reads the board through a view it computes
for itself; by two passive observers, one that pulls the board and one
that receives it pushed; and by a browser, in which a player and a
spectator are both written in JavaScript and reach the game over the
network. Six clients over one \texttt{Well} --- some in C\#, one a
PowerShell script, two in a browser --- and again the domain directory
does not change across any of them (Appendix A). Six is the number of
additions; counted by how differently the board is actually represented,
the labs hold two, a grid and a vector of column heights, and Appendix A
reports both numbers rather than the flattering one alone. What each
client adds is an input adapter, an output adapter, or both; none of
them is a change to the game.

The two experiments vary independent things. Experiment A changes where
the domain runs; Experiment B changes who observes it and how. Neither
touches the domain, and that independence is the result: its constancy
does not depend on where it is staged or on who is watching. The
identity §1 argued to be prior to the staging appears here as an
invariant the two variations cannot move --- and the build graph says
why it cannot be moved: nothing the domain refers to changed, because it
refers to no staging and to no framework in the first place.

How much these counts separate this arrangement from a competent
ports-and-adapters one is a question §9 answers with a baseline that was
built and counted, and the answer is less than the counts alone suggest:
a hexagon's rule model takes a new transport at no cost either. What the
counts establish here is the invariance itself, not a saving over an
alternative.

\subsection{3. What a Client Sees}\label{what-a-client-sees}

Experiment B added clients by adding adapters, and the word carries the
second half of the argument, so it is worth making precise. A client
never reads the domain directly. It reads through an adapter, and the
domain is the same behind every one of them.

What the domain emits is not a rendering --- but calling it a \emph{raw}
fact would overstate the case, and the difference needs a criterion
rather than an adjective. What is emitted is the set of occupied
interior cells, and that set is computed rather than stored: the domain
unions the settled pile's cells with the falling piece's and clips the
result to the frame's interior (\texttt{domain/Well.cs:322-331}). The
pile itself holds a set of positions (\texttt{domain/Pile.cs:26}). So
what leaves the domain is a derivation the domain performs, not a field
read off storage --- and the domain project, as §2 established, declares
no dependency on the framework or on any staging: it has no notion of a
screen, a renderer, or a client. The substrate pushes that set outward
on each move (\texttt{actor/TetrisActor.cs:43-47}), forwarding it rather
than projecting it further.

Derivation, then, cannot be the line, since the automated player's
column profile is a derivation too. The criterion this paper uses is the
one the series applies elsewhere, and it is a single question a reader
can put to any domain's source: \textbf{does a rule of the domain depend
on the notion?} If some operation is stated in those terms, the notion
is the domain's own, and what is emitted in them is a fact it holds; if
no operation mentions it and only a consumer needs it, it is a view
computed for a reader. Applied here the two separate cleanly. Occupancy
by position is the vocabulary the rules are written in --- collision is
a membership test over cells, and what counts as a complete row is
counted in them --- so the emitted set is in the domain's own terms.
Column heights, the gaps between them, their bumpiness appear in no
operation of the \texttt{Well}; nothing in the game is decided by any of
them, and only the automated player wants them. That is why the one is
emitted and the other computed outside.

What the criterion buys is that the line is inspectable rather than
asserted: a third party can read the domain and check which notions its
rules turn on. What it does not buy is a line that is sharp or
permanent. Which notions a domain owns is a modelling judgment --- Paper
8 called this the negotiable boundary of the two it drew --- and a
\texttt{Well} that grew an operation deciding something by column height
would move it, legitimately. The claim here is not that the boundary is
objective, only that it is auditable, and that in this domain the audit
comes out as this section describes.

Each client turns the fact into something it can use, and the turning is
the client's work, not the domain's. The human's on-screen grid is
produced by a renderer that lays the bitmap out as squares
(\texttt{actor/BoardRenderer.cs:19-48}). The browser produces the same
grid a second time, independently, in JavaScript --- its code is
commented as mirroring that renderer (\texttt{web/Program.cs:169-178})
--- so the terminal grid and the browser grid are two adapters over one
frame. The automated player needs something else: it reasons poorly over
a bitmap, so it lifts the same frame into a column-height profile --- a
skyline, the gaps and wells between columns, a few aggregate figures
(\texttt{tools/pile-scan.ps1:54-83}) --- a form suited to deciding where
a piece should go. None of these is the board's own appearance; nothing
in the domain decides anything by a picture. It emits a set of cells in
its own terms, and each client projects that through an adapter chosen
by the host that runs it.

One fact about how that adapter came to exist belongs in the argument
rather than only in the disclosure at the end of this paper, because it
is the sharpest evidence available for the claim of this section. The
automated player is an instance of a large language model, and the
column-height view it reads the board through was written \emph{by that
instance} during the lab, with the author's permission, to suit the form
in which it reasons (\texttt{tools/pile-scan.ps1}). Nobody translated
the board for it. It was handed the same emitted fact every other client
gets, and it authored its own projection of that fact into the shape it
could use. A projection belonging to the party that reads it is, in this
one case, not an interpretation of the architecture but an observed
event: the reader was a different agent from the domain's author, and it
built the adapter itself. How that came about, and what it means for
weighing the evidence, is disclosed in the acknowledgments.

This is where the on-screen grid loses its privilege. It is natural to
treat the pixels a human sees as the real board and the automated
player's vector as a derived abstraction of it. The system does not:
both are adapters over the same emitted fact, and neither is closer to
the domain than the other. The grid feels primary only because it is the
one a person happens to read; the domain has no such preference, because
it produces neither --- it produces the fact both are made from. The
automated player makes this hard to miss precisely because its adapter
looks nothing like a human's. Faced with a client that reads the board
as a vector of column heights, one notices that a human, too, only ever
reads it through a rendering, and had simply stopped noticing.

This arrangement has a much older statement in the database literature,
and the paper is better for saying so. Logical data independence ---
that an application should be insulated from the structure in which data
is held, and served instead by views derived from it --- was part of the
relational model from the beginning (Codd, 1970), and the three-schema
architectures that followed made the separation of a stored schema from
per-application views explicit. ``The domain emits facts and each client
projects them'' is that separation, reached from the side of the domain.
What differs is who owns the projection. Under data independence the
mediating system owns the schema and derives the views: the projection
is the database's work, performed on the application's behalf. Here
there is no mediating schema authority at all --- the domain emits its
own facts in its own vocabulary, and each client's projection is
authored outside both the domain and any central schema, by the staging
that binds it. The insulation is the same; the party doing the
insulating is not.

This connects directly to the boundary Paper 8 drew. There, what an
output \emph{is} --- its projection --- was shown to belong to the actor
that authors it, not to the domain that holds the material; a domain
object that renders itself reaches past what it knows into a decision it
does not own. Here the same boundary is seen from the client's side: the
projection is authored per client, outside the domain, and the domain
that would have rendered itself instead emits a fact and lets each
client author its own view. The adapters are the input sources and
output sinks of the earlier papers --- a keyboard or a pipe on the input
side; a terminal grid, a numeric vector, or an HTML canvas on the output
side --- bound to the actor from outside and exchanged without the
domain's knowledge. Whether a client is served by pulling or by being
pushed to is a choice of the same kind, made in the same place, and the
two modes have long been catalogued in the publish/subscribe literature
(Eugster, Felber, Guerraoui, \& Kermarrec, 2003). That the domain does
not know which adapter is attached is the same fact, seen once more,
that let the stage and the client vary in §2 without moving it.

\subsection{4. Actor and Audience}\label{actor-and-audience}

Section 3 answered \emph{through what} a client reads the domain: an
adapter, never the raw domain, and no adapter privileged. A second
question sits beside it --- what, in time, a client reads --- and the
answer is narrower than the temporal reading of it invites. Acting and
watching are distinct roles and what each may have differs; that is a
separation of roles, not a law about time. The distinction matters,
because the temporal version of the claim is either trivial or false.

Take the obvious shortcut first, since it is the one that has to be
dismissed with care rather than in passing. Each mutating command could
return the board it produced, and the client that issued the move would
have its picture at once. There is nothing incoherent in that. It is
\emph{read-your-writes}, one of the named session guarantees (Terry et
al., 1994); it is usually what a commanding client wants; and it is
often chosen for good reasons of latency. What such a response cannot do
is serve a client that issued no command, and that limit is close to
definitional --- no request, no response to be served from. So what
follows is a design corollary rather than a logical necessity: a system
whose only path to a view is the command's response has no path at all
for an observer that does not command, and every client in §2 and §3
that was not driving the game is exactly such an observer.

This is not a new mechanism, and --- the point deserves to be conceded
rather than discovered --- the temporal reading of it is not new either.
The field already separates the two operations: command-query
separation, and command-query responsibility segregation after it, hold
that a view is not served from the write (Meyer, 1988; Young, 2010), and
event sourcing builds a view as a projection over recorded acts rather
than a snapshot of current state (Fowler, 2005). And that what an
external consumer holds is the past was argued two decades ago in the
vocabulary of data: data outside a service lies beyond the reach of that
service's transactions and is \emph{temporally disconnected} from it, so
what has left is what was and not what is (Helland, 2005). Nor is the
structure beneath that new: in distributed systems it is the causal
past. Lamport (1978) showed that the events available to an observer
form a partial order --- what can have reached it is what happened
before it --- so an observer's knowledge is constitutively of the past,
and a ``present'' shared across separated parties is not a thing that
can be observed at all. This section therefore does not claim the
past-tense view as a finding; it sits inside a structure distributed
systems have had since 1978, and the data literature stated for
consumers two decades ago.

Nor is the present unavailable to an audience, and this paper's own labs
say so: one of the six clients polls the board and is handed the current
state (Lab C). An audience can see the present. What differs by role is
narrower --- a live read taken \emph{in order to act} belongs to the
party acting. The game's host queries the current board, is a piece
falling, so as to choose its next move, and that reading is part of the
act; an observer polling the same board is not deciding anything, and is
being shown a configuration. The difference is one of standing, not of
timing, and in that form the claim is not this paper's to make: Paper 4
gave the actor its voice, and Paper 8 established who is entitled to say
what an output is. This section applies their result to the seats in a
theatre rather than extending it.

Naming the two apart is where this paper needs the one before it. An
observer given only the present is given a snapshot --- the
configuration as it stands --- from which, to learn how that
configuration arose, it must reconstruct the history; that
reconstruction is the defect Paper 8 traced. An observer given the past
is given the account itself, the sequence of acts the domain performed,
which Paper 8 named \emph{testimony}: knowledge held on the word of the
one that lived it, not rebuilt from stills. The record-reading view of
this paper is therefore not merely consistent with Paper 8 --- it is
Paper 8's resolution made operational. To see a domain by reading its
record is to receive its narrative as testimony, and it is the only way
an audience comes to know \emph{how} the state arose. What polling
supplies is the configuration; what the record supplies is the account.

This returns to what the paper is about. Because what an audience is
\emph{told} is the recorded narrative rather than a state read off the
moment, the thing told is a standalone object: it exists in the record,
the same for every audience, independent of who reads it and when. And
it is the same for every audience for the same reason the board was ---
one domain, referenced by every staging, produced it. The invariant §5
names --- the domain as the common ancestor of its stagings --- is what
makes the narrative one narrative, readable the same from every seat.
What remains the same when the staging changes is, in the end, what
there was to see.

\subsection{5. Identity}\label{identity}

The two experiments leave a single structural property to state, and it
is statable without interpretation: a domain was run on five stagings
and read by six clients --- between them, in two genuinely different
representations --- and through all of it the domain referred to no
staging, every staging referred to it, and its content did not change.
In the dependency graph the domain is a sink; in the reversed graph it
is the root from which every staging is reachable, which is what §2
meant in calling it their common ancestor.

What that property licenses one to say is a second question, and the
word this paper uses for it invites more metaphysics than the evidence
supports. \emph{Identity} is offered as a reading of the property above
and not as a further finding: a thing that holds still under that much
variation is the sort of thing an entitlement can be attributed to,
which is what the rest of the series needs of it.

Two cautions follow, and they mark exactly where measurement stops. The
first concerns the word \emph{position}. A position in a graph
individuates nothing --- a different domain, or this one rewritten,
would occupy the same one --- so the property measured here cannot be
what makes this domain the domain it is; nor is it even the graph's only
sink, since the language's own library is another. The second is that
the property is evidence about a \emph{relation} and not a criterion of
identity: it says that what the stagings were built against is one
thing, and that its content owes them nothing. Calling that an identity
is a reading of the evidence rather than a part of it, and the two are
better kept apart than run together, because the measurement is
checkable and the reading is not. What this paper measures is
independence and priority. \emph{Identity} is the word for what
independence and priority license one to speak of, and the title should
be read in that light.

One echo in that title is worth disowning before a reader supplies it.
\emph{Identity precedes staging} has the shape of Sartre's
\emph{existence precedes essence}, and inverts it --- an essence prior
to any concrete appearance is close to what that formula was written to
deny. The resemblance is accidental and the claim is not the
philosophical one. Nothing here is an essence: what precedes a staging
is a position in a dependency graph and a source directory that did not
change, both of which can be read off a built artifact by someone who
has never heard of the phrase. A reader who hears the echo should hear
it as coincidence of form, not as an argument being smuggled.

Read that way, \emph{precede} is not a claim about time or ontology but
the direction of the dependence. A staging is built against a domain
that already exists and does not know of it --- the way a library is
compiled before, and in ignorance of, the programs that will call it.
The domain does not \emph{survive} its stagings; surviving would make
the stagings the primary thing and grant the domain the modest virtue of
enduring them. The relation runs the other way: the stagings were built
because the domain was there to be built against. Nothing the domain
refers to changed across the experiments because it refers to no staging
at all --- so its content was never at risk from a change downstream of
it.

This has a second face, less formal than the graph and worth stating
because it is what makes the identity legible to a person rather than
only to a build tool. A domain with an identity is one whose parts can
be understood as stable roles. To know what the board of the game
\emph{is} --- that it holds a pile of settled cells, spawns and moves a
falling piece, clears full rows, ends when a new piece cannot enter ---
is to know, structurally, what it does and does not do, what may be
asked of it, and what it will never answer. This is not a psychological
claim about how people understand stories; it is the observation that a
role with a fixed identity fixes its own boundary, and that boundary is
what domain-driven design reached for with its aggregates and bounded
contexts --- arrived at here by asking what stays the same, rather than
by drawing a diagram of layers. The invariance measured in §2 and the
legibility described here are the same thing seen twice: what holds
still across every staging is what there was to understand in the first
place.

\subsection{6. Recognition}\label{recognition}

If a domain keeps its identity across every staging, then whatever it
did keeps its identity too. The sequence of a domain's own acts --- a
piece spawned, moved, dropped; a row cleared --- is recorded in its
journal as the domain performed them, and because the domain is the same
on every stage, that record reads the same from every one. This is the
first, plain consequence of §5: the account of what happened is
available directly from the record the acts were written to, whole
wherever that record is kept, and does not have to be reassembled from
fragments held elsewhere.

That availability is what lets a client do more than watch. Given the
sequence of acts, a reader can recognize in it the routine the acts
compose --- that these three moves and a drop were the placement of one
piece, that this run of placements filled a row. The term needs pinning
down, since the rest of the section leans on it, and it is not this
paper's to define: a \emph{routine}, in the sense Paper 3 gave it, is a
pattern over journal entries together with the correlation that binds
them into one trajectory --- a reaction seeks an opening entry and then
a closing one, and what it matches, between them, is the routine. So
recognition is an operation with a definition and an implementation, not
a figure of speech, and it reads the record as it stands rather than
inferring a history from snapshots.

Recognition in this paper is that operation applied to a domain whose
identity holds: because the acts are the same across stagings, the
routine they compose is recognizable across stagings too. That step has
been measured (Appendix A, Lab H). One reaction, seeking a spawn and
then the drop that ends it, recognizes the same two placements in the
same order over three different stagings: the journal of a single
process, the journal a warm server keeps while driven over a named pipe,
and the journals of the three containers --- read inside each node,
against its own store. Within that third staging the three nodes'
records are byte-identical. \emph{Across} the three stagings what
matches is the sequence of acts --- same count, same order, same shape,
verb for verb --- while the entry identifiers differ by a constant, the
cluster writing three idempotent seeding acts where one process writes
one. The domain's diff for the whole exercise is empty. So the claim is
a result rather than an inference, on a short trajectory: two
placements, not two hundred.

That pair of results is worth reading together, because what moved and
what did not are not the same kind of thing. The identifiers moved: the
record's bookkeeping is a fact about where the domain ran. The acts did
not: a spawn on a console and a spawn inside a container are one
operation with the same parameters, and each leaves behind an entry of
the same kind. So what held still across the three stagings was not the
record's representation, which shifted, but what the act was and what
performing it left behind. A staging changes who is present when a
domain acts and changes nothing about what the acting does. That is the
sense in which §5's identity reaches past the source text --- the
invariance is not only in what the domain is written as but in what its
verbs accomplish when performed --- and it is Paper 4's account of an
actor's speech that makes this a claim rather than another way of
stating the empty diff.

Three things the lab found are more interesting than the confirmation,
and §6 is narrower for them. The first is that this record offers no
handle to correlate on: a spawn names a piece \emph{type} and a drop
names nothing at all, so the trajectory between them is tied by order
alone and not by any identity the acts carry. The second is worse, and
it answers part of the question this section leaves standing --- whether
a reading can be wrong. It can, and silently. Landing is not an act
here: the board lands a piece inside a tick or a drop, privately, so a
piece that comes to rest under gravity leaves its opening unclosed, and
the \emph{next} piece's drop closes it. The count comes out right and
the correlation comes out wrong, with nothing to signal the difference.
The third is that entry identifiers are not invariant across stagings
--- the three-node staging writes three seeding acts where one process
writes one, so everything after shifts --- and the closing identifier is
the only handle a match hands its reader. The routine recognized is the
same; the name by which a reader could refer to \emph{this instance} of
it is not.

The second of those generalizes beyond this lab, and it is the same
shape as the finding of Lab G. A routine is recognizable only in the
moments of it that were \emph{acts}. Where a moment was a private
consequence --- a landing folded into a tick, a pile absorbing cells ---
the record holds no entry, and no pattern can seek what was never
recorded. Recognition is bounded not by the reader's ingenuity but by
what the domain performed out loud.

That bound exposes a gap in this paper, and it is better named here than
left for a reader to find. Section 3 offered a criterion for the
\emph{vocabulary} of a record --- does a rule of the domain depend on
the notion? --- which is auditable against the domain's own operations.
There is no companion criterion for its \emph{granularity}: nothing in
this paper says which moments must be acts. Yet granularity is what
recognition rests on, so the fidelity of a record as an \emph{account}
depends on a decision the paper leaves unguided while the fidelity of
its \emph{vocabulary} has a test. And the failure mode is the worst
kind, the one Lab H found: an absent act yields not an error but a
smaller coherent story, matched to the wrong closing entry, with the
count coming out right.

No criterion is offered, and the honest statement is that this paper
does not have one. What its labs supply is three cases any candidate
would have to accommodate, reported as three cases and not as a rule.
Absorbing a landed piece had to become an act when the board was cut in
two, because two parties had to speak and the pile's own record needed
an act of its own (§8). A landing had to be an act for a reader that
\emph{reads} the record, since Lab H's reaction could not seek what was
never written. And a score and a difficulty level did not have to be
acts at all, being folds over acts already recorded (Lab I). Together
they point at what a \emph{reader} of a record needs to identify, as
against what a \emph{re-performer} of it needs --- but three cases are
not a criterion, and offering the hint as one would be the overreach §5
declines elsewhere.

It is worth locating the gap, since it is not peculiar to this
arrangement. Choosing which moments a model records is a modelling
decision of the same kind as drawing an aggregate's boundary, for which
the domain-driven design literature offers heuristics and worked
judgment rather than a test (Evans, 2003). What this arrangement changes
is only when the mistake surfaces: a boundary drawn badly resists the
code being written against it, whereas a moment that should have been an
act surfaces much later, in a reader that recognizes the wrong thing and
reports it to no one. That is an open problem this paper raises and does
not close, and it is worth being exact about what kind of question it
is. Which moments must be acts is not a question about representation
--- about how faithfully a record renders what happened --- but about
what the domain is \emph{doing} when it performs one. This paper has no
account of that.

One consequence of it belongs beside §8.2 rather than here alone,
because it is the same cost seen from the other end. Growth in this
arrangement is additive when the new thing is a value or a rule over
acts already recorded --- that is what Lab I measured. Discovering that
a \emph{moment} should have been an act is the other case, and it is the
expensive one, because every record already written lacks that act and
no projection over those records can supply it. What the domain would
have to do is not add an operation but reinterpret a history that does
not contain the distinction, which is exactly the schema-evolution cost
§8.2 concedes to the event-sourcing literature (Overeem, Spoor, \&
Jansen, 2017; Overeem, Spoor, Jansen, \& Brinkkemper, 2021). So the
missing criterion is not only an intellectual gap: it is the decision
whose mistakes are dearest to correct, and that is the strongest reason
to want one.

This is where the paper meets the one before it. Paper 8 followed an
output to the party that receives it and showed that, within a trust
boundary, the receiver comes to know the state as \emph{testimony} ---
an account it is told rather than one it reconstructs. But an account
received is not yet a routine recognized: being told the sequence of
acts is prior to seeing, in that sequence, what was done. This paper
takes the next step, and reaches it as a consequence rather than a new
claim --- an account whose subject keeps its identity can be recognized
as the same routine wherever it is told.

And the recognition is the client's, performed through the same adapter
§3 described. The board does not announce ``a row was cleared''; it
records the acts, and each client recognizes the routine through the
view it holds. The automated player reads its column-height vector and
recognizes a well being filled; a person reads the grid and recognizes
the same thing; neither recognition is the domain's, and neither is
privileged over the other. What both recognize is one routine, because
beneath both adapters is one domain with one identity --- which is where
the paper began. That no recognition is privileged does not settle
whether recognition is itself governed: who is entitled to read a
routine, and whether a reading of one can be wrong, is a question this
paper reaches and leaves standing.

\subsection{7. Where the Narrative
Lives}\label{where-the-narrative-lives}

Section 6's availability has a counterpart, and it wants naming
carefully. When what a system does is spread across autonomous services,
each with its own store, the sequence of what happened is in none of
them. It has to be assembled afterwards, by correlating records that
were produced separately --- which is part of what distributed tracing
and correlation identifiers are for. That work has a developed
literature: Google's Dapper (Sigelman et al., 2010, a technical report
rather than a refereed paper), and the refereed line that followed it
(Mace, Roelke, \& Fonseca, 2015; Sambasivan et al., 2016; Kaldor et al.,
2017). None of that is treated here as a deficiency, and nothing in this
paper says a system should do without it. The point concerns
\emph{completeness at a reader}: whether some single record holds the
whole account, or whether the whole exists nowhere and must be joined
together out of fragments.

Put that way, the section survives its own Lab B, which would otherwise
contradict it --- three machines keep three journals, converging by
replication and by a catch-up path, so the narrative is in three places
and reading it does depend on a mechanism. What makes that unlike
correlation is not that there is one copy but that the three are copies
\emph{of one record}: each holds the same acts, each alone holds the
whole game, and a reader at any node joins nothing to anything. So the
honest form of the claim is \emph{one logical narrative, replicated,
with convergence the framework guarantees} --- weaker than ``one
place'', and what the labs support. It is scoped, too, to one actor's
own first-person account (Paper 4): a system of many actors has many
accounts, and relating them is a question this paper does not take up
(§8). The model does not remove the need to know what a system did; it
changes where that knowledge already sits.

\subsection{8. Limits, and Two Observations Nobody Was Looking
For}\label{limits-and-two-observations-nobody-was-looking-for}

The claim is bounded, and the boundaries are worth drawing, the more so
because the result is easy to overstate into a promise it does not make.

\subsubsection{8.1 What the claim covers}\label{what-the-claim-covers}

It holds only where there is something for it to hold of. A domain with
a genuine identity, staged more than one way, is where the invariance
has teeth; a program too small to have a domain distinct from its
deployment, or one that will only ever run one way, has no second
staging to be invariant across, and the apparatus is overhead. The claim
is not that every program's domain is separable from its deployment in
practice, but that where the two are genuinely distinct, the separation
is real, buildable, and measurable --- and that collapsing them, where
they were distinct, is what the ordinary practice does.

A second limit, and the more consequential one, follows from what the
demonstrated domain is. The \texttt{Well} is a pure computational core:
it holds rules and state, and its outputs are data. It keeps no external
store, provokes no effect in the world, reads no clock, authorizes
nothing, and depends on no other party --- it never has to \emph{cause}
anything. That is exactly the case in which ``the domain emits facts and
each client projects them'' suffices, because every consumer really is a
client: it wants to know, and nothing turns on whether it ever does.

What the example does not exercise is a domain that must bring about an
external effect and then reason about the result --- charge a card,
notify a partner, settle a payment. There the consumer is not a client
but an obligatory collaborator, and ``the domain emits facts and each
client projects them'' is not a description of that relation. Nothing
here was built that way, so nothing here measures it. Paper 4 (Rivera,
2026d) takes up cross-actor continuity and its \texttt{tell}, which is
where a reader with such a domain would look; whether it settles this
case is not a question these labs asked.

\subsubsection{8.2 Invariance under staging, not under
evolution}\label{invariance-under-staging-not-under-evolution}

\textbf{Identity here is invariant under staging, not under evolution.}
The claim is that a \emph{change of staging} does not change the domain
--- not that the domain is frozen. Requirements grow, boundaries are
redrawn, and verbs are retired; a domain that does any of those is still
a domain no staging can call for a change to. Two measurements bound
that, one in each direction.

\textbf{The domain grew.} A score and a difficulty level were added to
the board --- the authority case, a value any client could compute and
which therefore has to be single-valued, and a rule over acts already
recorded (Appendix A, Lab I). Nothing was deleted but doc comments, no
signature changed and no verb removed, shown by diffing a sorted
inventory of the board's members. What the lab supplies that this paper
had not measured is the converse of §2: when the domain grew, all twelve
host projects kept running with no edit, checked by running them rather
than compiling them. The cost of \emph{adopting} the new capability then
fell only where it was wanted --- nothing for four of them, one line for
five, four lines for each of the two browser hosts, and thirty-two for
the input host, the one place where the level actually changes what
happens.

\textbf{The domain was re-cut.} One role was modelled too large and
became two, by authoring the two and then reading the original's
recorded acts --- the same read rehydration performs --- and driving
each new role to perform its own, so that each ends holding a record in
its own voice (Lab G). The original's journal is not cut or rewritten;
it is read as the account of what happened, and kept. Eleven of the
twelve host projects were untouched while the domain divided beneath
them, and the twelfth needed one line.

Where such a boundary may be cut has a criterion, and §8.4 relies on it:
\textbf{a boundary is admissible where no decision on either side reads
state the other party mutates concurrently.} The cut of Lab G passes it
because the board's phases are disjoint --- the pile is immutable while
a piece is falling, and the piece is gone once it has been absorbed ---
so neither role ever reads the other's state mid-change. The criterion
is about concurrency and not about size, and it is what makes a
tight-looking split sound.

That second one could look like a threat: if the boundaries
\emph{within} a domain can be redrawn, is the identity being measured
merely that of the current cut? The distinction the paper turns on
answers it. A re-decomposition is a change to the domain, made by
whoever authors it for reasons of modelling; it is not a change of
staging, and no staging can call for one. What carries across is the
record --- the acts were performed and remain performed --- so the
append-only character of every record involved survives the redrawing
untouched.

\textbf{The shape of growth is additive, and one measurement supports
that rather than establishing it.} An operation added is a new act in
the record rather than a reshaping of the acts already there, so growth
accretes; Lab I is that shape observed once, on the authority half of
the change. The half that adds an \emph{act} rather than a value --- a
piece held and swapped, a move undone, a second player --- is unclimbed
here, and so is the retirement of a verb, which in a production system
is marked obsolete and kept for as long as any recorded act still
invokes it. Both are named as future work rather than claimed. Nor
should additive growth be read as a claim that evolving a recorded
history is free: empirical study of event-sourced systems finds
converting recorded events as their schemas change to be a real and
recurring cost (Overeem, Spoor, \& Jansen, 2017; Overeem, Spoor, Jansen,
\& Brinkkemper, 2021). Nothing measured here bears on it. The invariance
reported is across stagings of a domain, not across revisions of it.

\subsubsection{8.3 Threats to validity}\label{threats-to-validity}

One threat to the result's validity cannot be argued away, and it
belongs in the open. Every staging and every client counted here was
written by the same author as the domain, and a metric of the form
\emph{the domain did not change} is gameable by construction: any code
can be held constant if all variation is pushed outside it, and an empty
diff is equally consistent with a domain designed backwards from a known
set of stagings. What would discriminate between those is a staging
nobody anticipated, introduced by a party with no part in the design.

The record supplies some of that and not all of it, and the part it
supplies is checkable in the example's own history rather than asserted
here. The domain's last commit is dated 29 June 2026 and it has not been
modified since. Three stagings and clients were added on 23 July,
twenty-four days later --- a gesture client, a constrained-device
client, and the three-machine deployment --- in working sessions
separate from the one that authored the domain, and the two clients
appear in no earlier artifact. The domain's diff across all three is
empty. Two of the three are not labs of Appendix A, having been
withdrawn there for want of measured evidence about the clients
themselves; what is cited here is only their effect on the domain, which
was measured, and which is all this particular argument needs.

It remains weaker than an independent test. The author commissioned the
work, and cross-machine deployment had been named as a next step in a
host comment the day after the freeze. One structural fact bounds part
of the worry without dispelling it: the fence is kept by the compiler
and not by anyone's discipline, since the domain declares no dependency
and its verb-bearing types are internal (§2), so staging concerns could
not have been threaded quietly into the domain even had someone wished
to. What that does not bound is the other form of the objection --- that
a domain's \emph{surface} may have been chosen with the known stagings
in mind --- and no property of a build graph rules that out. Nor would
another client on this same domain settle it, which is worth saying
because it is the obvious next move and the wrong one: the objection is
about who authored \emph{this} domain, so what answers it is the
arrangement applied to a domain its author did not design for it. That
is a second worked domain, not a seventh client, and this paper does not
have one.

A further limit is now measured rather than suspected, and it narrows
the result: invariance across \emph{transports} is not distinctive of
this arrangement. A ports-and-adapters version of the same game, built
and counted (Appendix A, Lab F), also takes a new transport without
touching its domain --- twice over. So the zeros of Experiment A, taken
as counts per staging, do not separate this arrangement from a competent
ported one; what separated them in measurement was a new
\emph{capability}, a game outliving its process, which cost the ported
domain a change and cost this one nothing. The claim to carry forward is
about capabilities a record supplies, not about stagings as such. And
the distribution stagings --- the two coordinated actors and the three
machines --- were measured against no baseline at all, since a ported
design has no distribution story to port; those zeros stand uncontested
rather than confirmed.

That correction has a limit of its own, and it runs the other way. Every
other limit in this section narrows what the paper claims; this one
narrows the concession just made. The baseline is the author's own. It
was written after the fact, by someone who already knew which property
it had to exhibit, and written deliberately clean --- its eleven rule
files differ from the journaled domain's by one line each, the
namespace. What it establishes is therefore what ports and adapters
\emph{permit}: a rule model with no project references, no packages and
no framework named, unmoved by a new transport. It is not evidence that
such domains are what the pattern yields in practice, and nothing of
that kind is claimed here in either direction. So the comparison does
its job --- it shows this paper's per-staging counts are not
distinctive, and that correction stands --- while establishing nothing
about the population of ported systems, which would take a study of a
kind not attempted here. A reader inclined to read Lab F as showing that
the arrangement argued for is already common should note that it shows
only that a competent instance of the alternative can be built, by the
same author, for the purpose; which is the same standing this paper's
own artifact has, and no more.

Two narrower limits belong here as well. The demonstration is one domain
on one framework, and it is an existence proof, not a measurement: it
shows the separation is buildable and that its mechanical cost is a set
of edits not made --- a domain diff that stays empty --- not that the
separation is worth its apparatus at scale, or cheaper in production,
which are questions for other work. And the distributed staging of §2 is
bounded as that section stated: its peers meet through an out-of-band
bootstrap, its coordinator role does not rotate, its transport trusts a
certificate on first use, and its replication reaches agreement through
a catch-up path rather than an uninterrupted stream. Each bounds the
deployment; none touches the domain's invariance, which is the only
thing the staging was there to test.

One of those two observations bears on this paper's own claim and so
belongs here rather than with them. A fact a domain emits must be
derivable within a single actor's state, and that is a property of the
substrate constraining a modelling decision \emph{inside} the domain ---
influence running against the direction this paper measures. It declares
no dependency, imports nothing, and its diff stays empty, so §2's
measured claim is untouched. What it does mean is that \textbf{declaring
no dependency on the framework is not the same as being unshaped by it.}
The influence is narrow, and it falls on how a domain may be divided
rather than on what it says.

One axis is deliberately left aside, and it closes the list of limits.
This paper varies where a domain runs and who observes it, and holds the
domain invariant across both; it does not treat which actor talks to
which --- the addressing and topology of a system of many domains --- a
separate question with its own answer, not folded into this one.

\subsubsection{8.4 Two observations that bound the
interpretation}\label{two-observations-that-bound-the-interpretation}

During the experiments two observations emerged repeatedly. Neither is
part of the claim advanced here; they are reported because they bound
its interpretation and motivate later work.

The first is a hazard, and the evidence for it is the strongest in the
paper: \textbf{an incomplete record answers plausibly.} Three labs met
it independently, by three routes, each taking it for a local incident
--- a routine correlated to the wrong closing entry with its count still
right, a truncated replay returning a board internally consistent and in
one case exactly correct, a rehydration silently dropping every
zero-argument act (Labs H, I, G). A gap in a record yields not an error
but a smaller coherent story. That is unlike \emph{gray failure}, whose
defining feature is differential observability, the application's view
diverging from the observer's (Huang et al., 2017): here there is none,
since every reader reads the same consistent record and no second view
exists to disagree. What detects it is therefore not another observer
but a ground truth carried forward, which gives a rule cheap enough to
state for anyone measuring a journaled system, not this one alone ---
\textbf{a replay measurement must carry the state recorded at play time
and assert against it}, or it will pass while being wrong. Two
qualifications: it is a rule about measuring rather than running, and
the exposure is not peculiar to this arrangement, since any reader
reconstructing state from an append-only record has it. It also meets
the one row of §9's table of Waldo's four differences marked as not
addressed --- partial failure of a \emph{peer} is exercised nowhere
here, while partial failure of a \emph{record} is what these labs found,
and it is quiet in a way a dead peer is not.

The second is a constraint, and it is narrower than the first in a way
worth stating before it is read: it is a property of \emph{this
substrate's} emitting plane rather than of journaled domains in general,
and unlike the first it rests on an argument with checked premises
rather than on a measurement --- nothing about the loss was measured,
because the split actor of that lab exposes no sink, no formatter and no
output wiring, so there was no push channel there to lose (Lab G carries
both qualifications). On this substrate, then: \textbf{a fact a domain
emits must be derivable within a single actor's state.} A complete frame
was a join over the pile's cells and the falling piece's, computed
inside the domain; after the board is cut in two those live in two
actors, and a projection on the emitting plane reaches only its own
actor's state, so neither role can push a whole frame (Lab G carries the
premises and their anchors). Emitting a joint fact therefore bounds the
admissible decompositions of a domain --- a boundary may be cut wherever
§8.2's concurrency test permits, except across a fact the domain must
emit whole. What this does \emph{not} touch is §3 or Lab D, whose counts
stand as taken, nor observation of a divided domain: a party reading
both records and joining them is an ordinary adapter, which is what §3
argues for. That reader was not built and nothing is claimed for it.
What closes is the push path.

\subsection{9. Related Work}\label{related-work}

The separation of a domain from the machinery that runs it is not a new
aspiration; reading it as a \emph{measured invariance} rather than a
design prescription is where this paper sits among its neighbours.

\subsubsection{9.1 What kind of claim the concessions are
against}\label{what-kind-of-claim-the-concessions-are-against}

What follows concedes each of this arrangement's mechanisms to prior
work, and every concession is meant literally. But prior art on the
parts of a configuration is not prior art on the configuration, and that
should be clear before the concessions begin. An architecture in the
foundational account is three things and not one --- its elements, the
\emph{form} that constrains them, and the rationale for that form ---
where form is the properties of and relationships among the elements
(Perry \& Wolf, 1992). Same elements under a different form, different
architecture. A form is what is claimed here: a domain that declares
nothing on either side, whose acts are recorded by a substrate it does
not name, staged by parties that name it. Each element of that is old
and this section says where each comes from.

The form has four constituents, and because the contribution is the form
and not the parts, they are set out together here rather than left for a
reader to gather from the places each is mentioned.

{\def\LTcaptype{none} % do not increment counter
\begin{longtable}[]{@{}
  >{\raggedright\arraybackslash}p{(\linewidth - 8\tabcolsep) * \real{0.2000}}
  >{\raggedright\arraybackslash}p{(\linewidth - 8\tabcolsep) * \real{0.2000}}
  >{\raggedright\arraybackslash}p{(\linewidth - 8\tabcolsep) * \real{0.2000}}
  >{\raggedright\arraybackslash}p{(\linewidth - 8\tabcolsep) * \real{0.2000}}
  >{\raggedright\arraybackslash}p{(\linewidth - 8\tabcolsep) * \real{0.2000}}@{}}
\toprule\noalign{}
\begin{minipage}[b]{\linewidth}\raggedright
\end{minipage} & \begin{minipage}[b]{\linewidth}\raggedright
Constituent
\end{minipage} & \begin{minipage}[b]{\linewidth}\raggedright
Credited to
\end{minipage} & \begin{minipage}[b]{\linewidth}\raggedright
In the baseline?
\end{minipage} & \begin{minipage}[b]{\linewidth}\raggedright
What it accounts for
\end{minipage} \\
\midrule\noalign{}
\endhead
\bottomrule\noalign{}
\endlastfoot
i & a rule model that declares no dependency & information hiding;
functional core, imperative shell & \textbf{yes} --- measured tie &
nothing distinctive \\
ii & entered by calls on its own operations, calling nothing outward &
inversion of control & \textbf{yes} & nothing distinctive \\
iii & declares nothing about what becomes of what it emits & implicit
invocation & no & the empty publicly callable surface \\
iv & its acts are recorded as it performs them & system prevalence;
event sourcing & no & history and durability at no cost to the domain \\
\end{longtable}
}

\textbf{Table 1.} The four constituents of the arrangement, what each is
owed to, whether the ported baseline of Lab F has it, and which measured
difference each accounts for. Every constituent is old; the claim is the
set.

Whether the \emph{form} is old was settled by measurement rather than by
assertion in a related-work section. The baseline holds the first two
and lacks the last two, and each surviving difference traces to one of
the two it lacks. That the domain has no publicly callable surface
follows from (iii): a port must be visible outside the domain that
declares it, and where nothing is declared there is nothing to expose.
That history and durability cost the domain nothing follows from (iv)
and not from (iii) --- a substrate could carry a domain's emissions to a
sink without recording any of them, and then a restart would find
nothing to replay. The record is what makes the second difference, and
it is named as a constituent for that reason.

\subsubsection{9.2 Work that isolates a
domain}\label{work-that-isolates-a-domain}

Most of that work can be compared on one property rather than mechanism
by mechanism, and comparing it that way is what the construct above
buys. \textbf{Existing work isolates a domain. None of it asks whether
the domain is finished.} Isolation is a relation between a domain and
its surroundings, and every line below achieves some version of it;
completeness is a property of the domain alone, and no line below
measures it.

{\def\LTcaptype{none} % do not increment counter
\begin{longtable}[]{@{}
  >{\raggedright\arraybackslash}p{(\linewidth - 4\tabcolsep) * \real{0.3333}}
  >{\raggedright\arraybackslash}p{(\linewidth - 4\tabcolsep) * \real{0.3333}}
  >{\raggedright\arraybackslash}p{(\linewidth - 4\tabcolsep) * \real{0.3333}}@{}}
\toprule\noalign{}
\begin{minipage}[b]{\linewidth}\raggedright
Prior work
\end{minipage} & \begin{minipage}[b]{\linewidth}\raggedright
How it isolates a domain
\end{minipage} & \begin{minipage}[b]{\linewidth}\raggedright
Asks whether the domain is complete
\end{minipage} \\
\midrule\noalign{}
\endhead
\bottomrule\noalign{}
\endlastfoot
information hiding (Parnas, 1972) & encapsulate the decisions most
likely to change & no --- a principle to follow, not a property to
check \\
domain-driven design (Evans, 2003) & a domain layer kept clear of
infrastructure & no --- and its repository is an obligation the domain
declares \\
functional core, imperative shell (Bernhardt, 2012a, 2012b; Peyton Jones
\& Wadler, 1993) & purity: effects confined to a shell & no --- its
subject is testability, and it is conceded below \\
model-driven architecture (Kleppe, Warmer, \& Bast, 2003; Mellor \&
Balcer, 2002) & platform independence, by transformation & no --- what
runs is a derivative, not the model \\
twelve-factor (Wiggins, 2011; Kratzke \& Quint, 2017) & configuration
separated from code & no --- its subject is the deployment, not the
domain \\
software product lines (Clements \& Northrop, 2001; Apel et al., 2013) &
one codebase, a variant per configuration & no --- it manages intended
difference \\
ports and adapters (Cockburn, 2005; Martin, 2017) & declared ports,
adapters outside & no --- the ports \emph{are} the incompleteness \\
location transparency (Hewitt et al., 1973; Agha, 1986; Armstrong, 2003)
& an actor reachable without regard to where it runs & no --- its
subject is the address of a peer \\
command-query segregation and event sourcing (Meyer, 1988; Young, 2010;
Fowler, 2005) & a view not served from the write & no --- the mechanism
§4 stands on, not a claim about the domain \\
\end{longtable}
}

\textbf{Table 2.} Prior work on isolating a domain, against the question
this paper asks of one. Three rows are also concessions of mechanism and
are developed below: ports and adapters, because it is the comparison
the result must survive; functional core, because it reaches the same
zero test doubles; and the two that own this arrangement's input and
output edges. Two borrowings of vocabulary are acknowledged rather than
compared --- \emph{accidental} is Brooks's (1987), used here of the
staging while his distinction partitions difficulty rather than
identity, so nothing reported here is a silver bullet; and the theatre
is Laurel's (1991), whose subject is the audience's experience where
this paper takes only the distinction between a play and its staging.
Brooks's cut was later applied to software's own materials by Moseley
and Marks (2006), separating essential logic from the accidental
complexity that state and control flow impose on it --- made within a
program where this one falls between a domain and its stagings, and the
reason the concession to functional core below is owed at all: keeping
an essential part free of accidental machinery is an old goal with a
literature.

That tradition also states the isolation goal this paper measures, and
it is worth naming the one point where its own answer stops short, since
a reader who practises it will feel that point before any other.
Domain-driven design asks that the domain layer be kept clear of
infrastructure, and for persistence it offers the repository: the
interface declared \emph{in the domain layer}, the implementation
outside it (Evans, 2003). That is a compromise rather than an oversight,
and Lab F shows why it is forced --- reconstituting an aggregate
requires a way in, so a design of that shape must have the domain offer
something. In the vocabulary above, the repository is an obligation the
domain leaves open, and a domain-driven design domain is decoupled
rather than closed. This one offers no repository interface, no factory
for a mapper, and no restore constructor, because reconstitution is the
recorded acts being performed again through the domain's own operations:
the way in is the way it was always entered.

The scope of that difference should be stated narrowly, because it is
easy to inflate into a claim about the whole tradition. It concerns
\emph{isolation from infrastructure} and nothing else. Domain-driven
design is also a discipline of modelling --- a language shared with
domain experts, contexts bounded deliberately, a model refined against
the business --- and a Tetris board demonstrates none of that. Nor does
this paper reach the domains that tradition is usually brought in for,
which act on the world and then reason about what happened; §8.1
concedes them. What is claimed is one axis: on the isolation goal both
share, the arrangement here pays no repository, and the ported baseline
built for comparison shows what paying one costs.

\subsubsection{9.3 Ports and adapters, and what closure adds to
it}\label{ports-and-adapters-and-what-closure-adds-to-it}

Closest, and the comparison the result must survive, is hexagonal
architecture --- ports and adapters (Cockburn, 2005) --- and its
restatement as clean architecture, the domain at the centre with every
dependency pointing inward toward it (Martin, 2017). The comparison is
usually put as a question about staging invariance, and put that way it
is settled and uninteresting: a ported version of this same game was
built, given the same stagings one at a time, and counted, and two of
its transports cost its rule model nothing at all (Appendix A, Lab F).
Per staging the two arrangements are indistinguishable, and this paper
concedes it without reservation.

The difference that does separate them is not a count of edits, and it
is easier shown than argued. Set what each arrangement still requires of
someone else side by side, taking every row from Lab F or Lab E:

{\def\LTcaptype{none} % do not increment counter
\begin{longtable}[]{@{}
  >{\raggedright\arraybackslash}p{(\linewidth - 6\tabcolsep) * \real{0.2500}}
  >{\raggedright\arraybackslash}p{(\linewidth - 6\tabcolsep) * \real{0.2500}}
  >{\raggedright\arraybackslash}p{(\linewidth - 6\tabcolsep) * \real{0.2500}}
  >{\raggedright\arraybackslash}p{(\linewidth - 6\tabcolsep) * \real{0.2500}}@{}}
\toprule\noalign{}
\begin{minipage}[b]{\linewidth}\raggedright
\end{minipage} & \begin{minipage}[b]{\linewidth}\raggedright
What is measured
\end{minipage} & \begin{minipage}[b]{\linewidth}\raggedright
Ported baseline
\end{minipage} & \begin{minipage}[b]{\linewidth}\raggedright
Here
\end{minipage} \\
\midrule\noalign{}
\endhead
\bottomrule\noalign{}
\endlastfoot
\textbf{I} & \textbf{outward obligations the domain declares} &
\textbf{3 driven ports} --- board output, piece selection, state &
\textbf{0} \\
↳ & stand-ins its own tests require & 3, without which 20 of 64 tests do
not run & 0 \\
↳ & publicly callable domain operations & at least one per port & 0 ---
one type with no members \\
↳ & a state admissible from outside & \texttt{GameState}, arriving
through the state port & none; only acts \\
\textbf{II} & \textbf{a reconstitution surface of its own} & \textbf{a
restore constructor and a pile factory: 56 lines added, 5 removed} &
\textbf{none} --- replay re-enters by the operations the acts were
performed through \\
\end{longtable}
}

\textbf{Table 3.} What each arrangement still requires of someone else
--- \textbf{two independent measurements, not five.} The rows marked ↳
are consequences of the fact above them and are shown because each is
separately checkable, not because each is separate evidence: a stand-in
is a driven port seen from the test side; a port must be visible outside
the domain that declares it, so where nothing is declared there is
nothing to expose; and the state the baseline will accept arrives
through one of those same three ports, the one its own source calls
\emph{a third driven port, added by staging 4}. The quantity is
unfulfilled obligations, not lines of code.

The conjunction of those two measurements is what this paper calls
\textbf{closure}, and the word is doing no work the counts do not:
\emph{the domain declares no obligation outward, and its reconstitution
requires no surface of its own.} A port is an obligation the domain
itself declared and left for someone else, so the first fact is what
makes the three rows under it follow; and reproducing by the operations
the acts were performed through is what makes the second one hold, since
there is no event distinct from the act for a separate replay surface to
take. A domain with neither obligation is finished. Ports and adapters
\emph{decouples} a domain; this arrangement \emph{closes} one.

Two established senses of the word meet here, and both are better dealt
with than left for a reader to supply. One is favourable and is claimed
deliberately: in the language of modules and linking a component with no
unresolved external reference is \emph{closed}, and the closure of a set
of modules is what a linker computes when nothing outward remains to be
found. That is very nearly the property measured above, arrived at by
asking a different question, and the resemblance is not a coincidence to
be excused. The other is unfavourable and is denied. In the open/closed
principle, \emph{closed} means closed to modification (Meyer, 1988) ---
and this domain is not that. It grows: a score and a difficulty level
were added to it, additively and without a signature changing (§8.2, Lab
I). What is closed here is the set of obligations the domain leaves to
others, not the set of changes its author may make to it.

Three things bound that, and the first is the most likely to be misread.
Closure is not visible in the domain \emph{considered in isolation}: a
hexagon's driving port is the domain's own API, which nobody supplies
and nobody doubles, and a domain needing nothing outward declares no
driven port either --- the \texttt{Well} as bare rules over a board is
such a domain (§8.1). All five rows appear only once the domain is
\emph{deployed}, when its state must outlive a process and arrive on a
second machine. In the ported arrangement those became declared needs;
here they were never needs, the acts having been recorded as they were
performed.

The second bound is on the word. Calling the sum of those counts
\emph{closure} is a reading of them, as §5's \emph{identity} is a
reading of the structural property, and the paper keeps reading and
measurement apart here as it does there. The third is that history and
durability are one row of the table and not the claim; purity, conceded
below to functional core and imperative shell, buys freedom from effects
and closes none of these rows. The obvious objection --- that orthodox
practice persists an aggregate through an adapter-side mapper without
touching the rule model at all --- is answered by measurement rather
than by argument: the baseline attempted exactly that design and still
needed a seam authored inside its own rule model, for a reason its
source states and Lab F sets out.

One measured result runs the other way and belongs in this comparison
rather than in its appendix. The build graph does not separate the
arrangements: a ported rule model is equally free of project references,
packages, and framework names. So the testability argument above must be
narrowed to what it can bear. What cannot be exercised without a
stand-in is the ported \emph{application service} and its three driven
ports, without which twenty of its sixty-four tests do not run; the
ported \emph{rule model}, taken alone, is as testable and as
dependency-free as the journaled domain. The driven and driving sides
are not symmetric.

It is fair to press the point: if the substrate provides the mechanism,
has the port not simply moved into the framework --- universal, unnamed,
but a port all the same? Two answers, and the second is the one that can
be checked. A term stretched that far stops separating anything: a
domain declares neither the garbage collector, nor the scheduler, nor
the collections library it computes with, and calling each of them a
port empties the word. What \emph{port} earns its keep by naming is an
interface a component \textbf{declares and depends upon}, because that
is the thing with the consequences this paper counts --- an
implementation per staging, a stand-in per test. And measured there the
difference is not interpretive: the domain of §2 compiles and passes its
tests with the framework absent from its build graph, which a domain
that declares a driven port cannot do, since a port with neither
implementation nor double is not something a test can run against at
all.

\begin{verbatim}
   Ports and adapters:
       Adapter --calls--> [driving port = the domain's OWN api] --> Domain
       Domain --declares--> [driven port] <--implements-- Adapter
       (only the driven port is a dependency of the domain, and a domain
        that needs nothing outward declares none)

   Here:
       Domain        Staging --binds--> input source / sink / client
       (no declared dependency either way; the state that must outlive the
        process and reach a peer is the substrate's doing, not a port the
        domain asks for)
\end{verbatim}

The contract has not been inverted; it has moved off the domain's
boundary, and two places are involved rather than one. The
\emph{contract} is the substrate's: it defines what an output is, what
may command the domain, and how either is carried. The \emph{binding} is
the staging's --- which sink, which client, which transport --- and it
happens later than a reader is likely to assume, a destination being
attached to a live actor rather than fixed at build or deployment
(\texttt{Choreography/Theater/PerformanceV2.cs:415}). So the formulation
used throughout is that the contract is provided by the substrate and
bound in the staging, never that it moved into the staging, which would
credit the staging with a mechanism it only wires up. The same shape
holds for the transport carrying an actor's messages to another, the
domain naming no wire there either; that belongs to the axis this paper
sets aside (§8) and to Paper 4, and is mentioned so the output side is
not read as a special case. What the move costs is measured rather than
reasoned (Appendix A, Lab F): a ported arrangement pays no domain edit
per staging either, and its driving side needs no double, a hexagon
implementing its driving port rather than depending on it. Where it pays
is on the driven side --- three ports without stand-ins for which twenty
of its sixty-four tests do not run --- and in its publicly callable
surface, at least one operation per port against none here. That is the
observable trace of the move and not the point of it. The point is that
the boundary the contract sits on is no longer inside the domain, which
is what lets the domain be the common ancestor of §5: a port is a thread
from the domain back to the shape of its clients, a fragment of a
staging lodged inside it, and while it is there the identity is not
clean.

\subsubsection{9.4 Who owns the two edges}\label{who-owns-the-two-edges}

The geometry itself has a prior claimant, and it is the nearest
neighbour this paper has. Implicit invocation --- a component announces
an event, components registered elsewhere are invoked, and the announcer
names no recipient and declares no interface for one --- was formalized
as a design space three decades ago (Garlan \& Notkin, 1991) and studied
as the means by which independently built tools are integrated and
independently evolved (Sullivan \& Notkin, 1992); its later forms are
surveyed as publish/subscribe (Eugster et al., 2003). A domain that
emits facts under logical names and says nothing about who receives them
is doing that, and the coupling that remains --- agreement on the names
--- is what that literature calls nominal coupling. The geometry of §2
is therefore implicit invocation's geometry, and the honest statement is
that this paper did not discover it.

Three things distinguish what is claimed here, and all three are
narrower than the geometry. The first is the direction the vocabulary
faces. In implicit invocation the event vocabulary is an
\emph{integration} vocabulary: announcer and listeners agree on a set of
events in order to be integrated, and the names exist because something
will listen. The names a domain emits here are its own --- the
operations it performs and the facts they leave behind --- and they
would stand unchanged if nothing listened at all; the staging adapts to
them, never the reverse. The second is that implicit invocation is
offered as a mechanism, with a design space to explore and trade-offs to
weigh, whereas the claim here is an invariance measured on a built
system: the domain's diff across five stagings and six clients. The
mechanism makes that zero possible; it neither asserts nor measures it.
The third is the conclusion drawn. That literature concludes that
components may be integrated and evolved independently. This paper
concludes that the domain stands as the common ancestor of every staging
of it --- a claim about what a domain is, not about how the parts of a
system are connected.

The input edge has a prior claimant of the same kind. That a framework
calls the application's code rather than the reverse was named
\emph{inversion of control} when frameworks were first described as a
form of reuse (Johnson \& Foote, 1988), and it is the mechanism by which
a domain can be driven without declaring an interface for whatever
drives it: the substrate invokes the domain's operations, and the domain
calls nothing outward. Section 2's input-side zero rests on that
mechanism and the credit belongs there. What the term later came to name
in practice --- a container that wires declared dependencies together
(Fowler, 2004) --- is a different thing, and the distinction matters for
the reason Paper 8 gave: injection presupposes something to inject, and
where a domain declares no port there is nothing to wire. The mechanism
is inversion of control; the measurement, and the identity it reveals,
are what is added here.

One practitioner pattern has arrived at that same geometry from a
different motive, and it deserves an explicit concession. Functional
core, imperative shell places a dependency-free core inside a shell that
owns all input and output, and its best-known consequence is precisely
the count reported above: the core can be tested with no test doubles,
because there is nothing there to double (Bernhardt, 2012a, 2012b). That
zero is therefore not this paper's discovery, and it should not be
presented as one. The name belongs to practitioner literature, but the
discipline beneath it does not: confining effects to a boundary so that
the remainder can be reasoned about and exercised without them is the
monadic treatment of input and output, argued in the refereed record
three decades ago (Peyton Jones \& Wadler, 1993). What differs is the
claim the zero supports. The pattern buys its core's independence with
\emph{purity} --- the core is functions without effects and the shell
holds the mutation --- and it is a discipline about the internal
structure of a single program, whose shell is one imperative wrapper.
The domain here is not pure: it mutates. What is recorded of that
mutation is recorded by the substrate, and the domain never learns a
journal exists --- which is the measured claim of §2 and not a courtesy
of this comparison. So purity is neither available as the mechanism nor
needed as one: the independence is bought not by declining to mutate but
by keeping the account of the mutation someone else's work. And its
shell is not one wrapper but many stagings, one of them spread across
three machines. Same zero, different result: the pattern explains why a
\emph{test} needs no double; this explains why a \emph{staging} needs no
domain edit, and what that says about the domain's identity.

\subsubsection{9.5 The local/remote
objection}\label{the-localremote-objection}

Which raises the strongest objection available to this paper, and it is
an old one. Waldo, Wyant, Wollrath, and Kendall (1994) argued that the
distinction between local and remote cannot be papered over: latency,
memory access, concurrency, and partial failure are irreducible, and
systems that hide them fail precisely where reliability is demanded.
Experiment A --- a domain moved from one process to three machines with
no edit --- has the shape of the thing they warned against. The
objection is accepted here rather than rebutted. Nothing in this paper
claims those differences vanish, or that distribution is free. Nor does
it help to gesture at the incidents a distributed staging happens to
produce --- an out-of-band rendezvous, a connect-readiness race, a
certificate trusted on first use --- as though each answered one of
Waldo's four. It is more useful to set the four out and say, for each,
where it falls here and whether this paper does anything about it:

{\def\LTcaptype{none} % do not increment counter
\begin{longtable}[]{@{}
  >{\raggedright\arraybackslash}p{(\linewidth - 4\tabcolsep) * \real{0.3333}}
  >{\raggedright\arraybackslash}p{(\linewidth - 4\tabcolsep) * \real{0.3333}}
  >{\raggedright\arraybackslash}p{(\linewidth - 4\tabcolsep) * \real{0.3333}}@{}}
\toprule\noalign{}
\begin{minipage}[b]{\linewidth}\raggedright
Waldo's difference
\end{minipage} & \begin{minipage}[b]{\linewidth}\raggedright
Where it falls here
\end{minipage} & \begin{minipage}[b]{\linewidth}\raggedright
Status
\end{minipage} \\
\midrule\noalign{}
\endhead
\bottomrule\noalign{}
\endlastfoot
Latency & the transport, and the replication that crosses it & present;
nothing about it was measured, no timing was taken \\
Memory access & nothing is shared by reference at all: each node keeps
its own copy of the record and reconstructs its own state from it &
answered by the arrangement rather than borne by it \\
Concurrency & the connect-readiness race of Lab B; one writer per record
& present, and the race is reported \\
Partial failure & not exercised --- no peer was killed, no partition
induced (Lab B) & \textbf{not addressed} \\
\end{longtable}
}

\textbf{Table 4.} Waldo et al.'s four differences against what this
paper does with each. Two incidents a reader might expect to find here
belong to neither column: an out-of-band rendezvous is a bootstrap
question and trust-on-first-use a security one, and neither is among the
four.

Read that way the honest statement is narrower and more interesting than
the one it replaces. \emph{Memory access} is the case where it hurts
most, and it is not evaded: there are three copies of the board's state
in Lab B and no shared reference anywhere, which is not a workaround for
the absence of shared memory but an arrangement that never assumed it.
\emph{Partial failure} is where this paper is weakest, and Lab B says so
already. What holds is that the differences which do appear are borne by
the staging, and that none of them reached the \texttt{Well} --- a claim
narrower than location transparency's, and deliberately so: the
local/remote distinction is real and irreducible, and it is not a fact
about the domain.

\subsubsection{9.6 How the claim is
checked}\label{how-the-claim-is-checked}

The method by which the claim is checked has a literature of its own.
Software reflexion models compare a high-level model of a system against
the dependency graph extracted from its source, reporting where the two
converge, diverge, and fall silent (Murphy, Notkin, \& Sullivan, 1995,
2001). The conformance-checking work that followed made such comparisons
routine: languages in which permitted and forbidden dependencies are
declared and then checked against the code (Terra \& Valente, 2009),
comparative accounts of the available static techniques (Passos et al.,
2010), and tools --- ArchUnit, NDepend --- in which such rules are
asserted as ordinary unit tests. What motivates that whole literature is
architectural erosion: the drift of a built system away from the
architecture intended for it (Perry \& Wolf, 1992). Section 2's reading
of the build graph is a measurement of the same kind, and the
resemblance is what lets this paper say \emph{checkable} rather than
\emph{arguable} --- with erosion as the instructive contrast, since it
is what a dependency direction does when nobody is enforcing it, and the
empty diff is what the same direction looks like when it is a property
of the arrangement rather than a rule someone must police. The
difference lies in what is checked. Conformance checking takes an
intended architecture as its input and counts departures from it; there,
the direction of dependence is a rule imposed and then policed. Here the
direction is not a rule but the substance of the claim, and the quantity
measured is not a violation count but an invariance under deliberate
variation: change the staging, and see whether the domain moves. The
first is an audit; the second is an experiment.

\subsubsection{9.7 What §4 owes the data
literature}\label{what-4-owes-the-data-literature}

The nearest neighbour of §4 is a data-management literature, and it must
be credited carefully, because two distinct papers bear on it.
\emph{Immutability changes everything} argues for append-only computing:
observed facts recorded and kept, results derived from them on demand,
on the model of an accountant who corrects by adding an entry rather
than by erasing one (Helland, 2015); Kleppmann and Kreps (2015) turn
that into a way of composing systems around a log. The earlier paper is
the one that reaches §4's temporal claim. \emph{Data on the outside
versus data on the inside} holds that data outside a service is beyond
the scope of that service's transactions and temporally disconnected
from it (Helland, 2005) --- which is to say that what an external
consumer holds is always the past. Section 4 therefore does not discover
that; it arrives at it twenty years later, in the vocabulary of roles
rather than of services, and the honest statement is that the data
literature drew the conclusion first.

What survives the concession is narrower, and worth stating exactly.
Helland's argument concerns where data sits relative to a transaction
boundary and what may be assumed of it there. Section 4's concerns
\emph{offices}: that watching and acting are distinct roles; that the
present is available only in the acting one; and that an audience is
therefore constituted by the record rather than merely served from it.
Its operative consequence is a design corollary the data framing has no
occasion to state --- a command's response cannot serve a client that
issued no command, so a system whose only view path is that response has
none for an observer --- which follows not from a property of data but
from whose the response is. The temporality is Helland's, and the
structure beneath it Lamport's; the attribution of roles, and the
corollary it yields, are what this section adds. It does not claim more
than that, and §4 says why: returning state in a command's response is a
recognized session guarantee (Terry et al., 1994), not an error, and the
present is observable by an audience that polls for it.

\subsubsection{9.8 The series, and what remains
unclaimed}\label{the-series-and-what-remains-unclaimed}

Two of this series' own papers hold pieces of the present result. Paper
7 read the server as an \emph{accidental category} --- ephemeral under
this substrate, discharged at bootstrap and absent from the running
topology after --- and offered that as a lens rather than a result. The
consequence drawn here is that where a domain runs is then not
constitutive of it, which is one axis of the staging this paper varies.
Paper 8 located, from the side of output, the boundary between what a
domain holds and what an actor projects for a client, which is the
adapter boundary of §3 seen from within a single output. This paper
generalizes both: it varies the whole staging --- where a domain runs
and who observes it --- and names what stays fixed across all of it as
the domain's identity.

What remains unclaimed, after all of these, is the intersection ---
which is to say the form, in the sense this section opened with. Late
binding of a destination is old; platform independence has a name and a
methodology; location transparency has a production tradition; a
dependency-free core has a pattern; being driven without declaring an
interface for the driver has a name three decades old; and checking a
dependency graph against an intended model has a literature. Each of
those is conceded above, individually and without reservation. What none
of them states is that the contract may leave the domain on \emph{both}
sides at once, so that a domain declares nothing about what drives it or
about what it emits --- and that the property this leaves behind, the
domain standing as the common ancestor of every staging of it, is
measurable after the fact on a system already built.

A reader is entitled to ask why a list of conceded parts should leave
anything behind, and the answer is the one the section opened with
rather than a rhetorical one. The claim is about a form, and a form is
tested by building the alternative. The alternative was built, from
parts every one of which is credited above, and it produced neither of
the two differences that survive. Nothing here rests on the parts being
new. What it rests on is that assembling them the ordinary way does not
arrive here --- a claim about the arrangement, and one this paper
publishes a falsification test for: build a ported design that obtains
durability and a second machine without touching its rule model.

One such design is available and needs naming, since it is the first a
reader who has built one will reach for, and it is nearer than the
mapper of Lab F. An event-sourced hexagon declares an event-store port
and reconstitutes by replay rather than by mapping, so it needs neither
a restore constructor nor a factory for a pile, and the 56 lines this
paper counts would not be spent. What it does need is a
\emph{reproduction surface} distinct from its command surface: an
orthodox event-sourced aggregate exposes something of the shape of
\texttt{Apply(event)} and \texttt{LoadFromHistory(events)}, so the rule
model still grows a seam, only a different one. The residual claim
stated exactly is therefore not that a ported design must pay a mapper's
seam, but this: reproducing by the domain's \emph{own} operations
requires no reconstitution surface at all, because the substrate
re-invokes the very verbs the acts were performed through, and there is
no event distinct from the act for an \texttt{Apply} to take. That is
\emph{predicted} rather than measured --- no event-sourced baseline was
built, and building one is the lab that would close the test this paper
publishes.

\subsection{10. Conclusion}\label{conclusion}

This paper began where practice is settled. A program's domain and its
deployment are written together, so changing where a system runs, or
which client it serves, means changing the system. Set beside theatre,
where one script is staged in many venues without being rewritten, that
entanglement stops looking like a property of software and starts
looking like a confusion of a play with its staging.

The confusion has a measurable opposite. One domain --- a Tetris board
--- was run on a console, in a browser over a WebSocket, in a browser
over server-sent events, across two co-hosted actors over real TLS, and
across three separate machines in containers; and it was read by a
person at a keyboard, an automated player over a pipe, two passive
observers, and a browser drawing it in JavaScript. The number of changes
that variation forced on the domain is zero, and the number of test
doubles the domain's own tests require --- for what drives it or for
what it emits --- is also zero.

Behind those counts is a property of the domain rather than a
discipline, and the counts are not where it lies: a ported rule model is
equally free of references and packages, equally testable without
stand-ins, and equally unmoved by a new transport (§9), so those counts
distinguish a clean domain from an entangled one and not this
arrangement from that one. \textbf{Ports and adapters decouples a
domain; this arrangement closes one.} Section 9 puts that on two
measurements --- whether the domain declares any obligation outward, and
whether reconstituting it needs a surface of its own --- and the 61
changed lines a ported rule model spent to survive a restart are one row
of the second, not a separate result. The domain here declares no
interface for what drives it or for what it emits, so the contract
fixing what an output means, where it goes, and what may command it does
not sit on its boundary at all: the substrate provides it and the
staging binds it. That is not dependency inversion carried a step
further; it is the contract in a different place, and it is what leaves
every staging's references leading to the domain and none leading out
--- a sink in the dependency graph, the root of its reverse, which is
all the paper means in calling the domain their common ancestor. What is
checkable against a built artifact is that structural property together
with the empty diffs; reading it as an \emph{identity} is what §5 does,
on the terms §5 sets. Read so, \emph{precedes} names the direction of
the dependence. A staging is built against a domain that does not know
of it; the domain is not the thing that survives its stagings but the
thing that makes them possible.

Two consequences followed with no new machinery. Watching and acting are
distinct roles, and what each may have differs: a live read taken
\emph{in order to act} is part of the act and belongs to the party
acting, while a view served from a command's response can serve no
client that issued none. So an audience is constituted by the record
rather than served a state read off the moment --- not because the
present is closed to it, since one of the six clients polls the board
and is handed it, but because what an audience is told is an account and
not a configuration. And because the domain keeps its identity, the
account of what it did keeps its identity too, so the routine those acts
compose is recognizable from any stage, through any adapter. Neither is
a second claim; both are what the first one entails.

What can be carried away is an instrument rather than a prescription,
and the difference is worth being exact about, since this paper declares
itself the analytic kind of contribution. Nothing here establishes that
an architecture is better for holding the property, which would be a
judgment about what makes an architecture good. What the paper supplies
is a question that can be put to a system already built and answered by
measurement rather than opinion: does the domain inside it hold an
identity independent of where it is staged and to whom? Vary the
staging, and count what the domain had to change. The labs of Appendix A
are that question asked once, of one domain, with the answer zero.

Two things came out of those labs that nobody set out to find. Neither
is the result of this paper, and both are reported because they bound
how its result may be read --- the first by bounding what a reader of a
record can be told, the second by bounding how a domain that keeps one
may be cut. The first is a hazard, and the evidence for it is the
strongest in the paper: \textbf{an incomplete record answers plausibly.}
Three labs met it independently, by three routes, none investigating it
--- a routine correlated to the wrong closing entry with the count still
right, a truncated replay returning a board that was internally
consistent and in one case exactly correct, a rehydration silently
dropping every zero-argument act. A gap in a record yields not an error
but a smaller coherent story, and because every reader reads the same
record and it is consistent, no second view exists to disagree. The rule
that follows is cheap and is offered to anyone measuring a journaled
system, not just this one: carry the state recorded at play time and
assert a replay against it. A replay checked only for internal
consistency will pass while being wrong. The second is a constraint, and
it is stated in §8: a fact a domain emits must be derivable within a
single actor's state, so emitting a joint fact bounds how that domain
may be divided --- the one respect in which the substrate shapes the
domain rather than the other way round, and the more interesting for
being the exception. Whether an architecture \emph{ought} to have the
property --- whether the arrangement pays for itself, and against what
--- is exactly what §8 and Appendix A decline to claim, and a reader who
wants that will have to measure it.

Four questions the argument opened are left standing, each larger than
this paper. If an entitlement must be attributed to something that holds
still, then identity is prior to authority, and the series' earlier
questions about who may speak rest on this one. Who is \emph{entitled}
to read a routine, §6 reaches without settling --- though whether a
reading of one can be wrong it does settle, and the answer is that it
can, and silently, when a moment that mattered was never an act. Which
is the third question, and the one this paper is most exposed on: what
must be an act at all. A criterion is offered here for a record's
vocabulary and none for its granularity, and it is granularity that an
account's fidelity rests on (§6). And because a domain that declares no
dependency cannot reach for another, whatever composes two domains is
neither of them --- the question that anything treating a set of domains
as a repertoire would have to answer first. Closure is what this paper
hands to that question. A domain that still declares obligations is a
\emph{part} of a system, completed by whatever supplies them; one that
declares none is a whole, and a whole is the only kind of thing a
repertoire can hold. Whether closure is sufficient for that, or merely
necessary, is not settled here.

The audience changes. The stage changes. The play asks nothing of
either.

\subsection{Appendix A. Labs}\label{appendix-a.-labs}

The claim of §2 is a count, and this appendix reports the counts and
says where each can be checked. These labs are illustration rather than
evaluation: each shows that a separation the paper argues for is
\emph{built}, and what it costs the domain, which is nothing. In the
manner of the earlier systems papers, every result here is a count the
separation drives to zero, and the zeros are the claim. Two things are
not measured and are marked as such rather than mixed in: whether the
arrangement is cheaper or faster at scale, or better in production,
which §8 leaves open; and what the distribution zeros are zeros
\emph{against}, since the one alternative arrangement that was built ---
the ported baseline of Lab F --- has no distribution counterpart to set
beside them.

All of them run against one domain --- the \texttt{Well} and its
operations (\texttt{domain/}) --- reached by every host through a single
actor (\texttt{actor/TetrisActor.csproj:11}), which the example's own
comment calls an accidental shell (\texttt{actor/IGameHost.cs:14}). The
domain project declares no reference of its own (Lab E).

\textbf{Lab A --- the stage.} The same \texttt{Well} runs on a console;
in a browser with input and output carried over a WebSocket; in a
browser over HTTP with server-sent events; and across two StageManager
actors, joined once in memory and once over a real Kestrel TLS channel.
Five stagings, differing in process, transport, and wire format. A diff
of the domain directory from the first staging to the last is empty, and
the example builds with no errors. The domain's own test suite also
passes unchanged throughout, which is worth stating for what it is and
no more: a regression check showing the stagings did not break the
domain. Its coverage was not measured, so neither the suite's size nor
its passing is offered here as evidence that the domain is adequately
exercised --- a test count is not a coverage measurement, and nothing in
this appendix should be read as one.

\textbf{Lab B --- three machines.} Three Docker containers run the same
\texttt{Well} as three StageManager peers on a private bridge network:
one Director, two casts, joined over Kestrel TLS on a port never exposed
to the host, with coordination and replication crossing genuine
container-to-container TLS. A scripted driver on the Director plays the
game to a non-trivial board; the three nodes converge on the same
journal entry and the frame each pushes is byte-identical. Adding this
staging left both the domain and the actor untouched --- the diff of
each across the whole increment is empty, and the files added are a new
host, its Docker definitions, and notes.

Two details of that run belong here rather than in the body. The game
was played by a script rather than by a person, but it was played: the
\texttt{Well} genuinely mutated on one machine and the same state
reached the others. And what each node is handed is the domain's
\emph{assembly}, named through the anchor type --- the \texttt{Well}
itself is passed nowhere and constructed by no host, but built inside
each node by a recorded act, the seeding command the actor issues on
startup, so that a board exists on three machines without any of them
ever having held one.

The weight this lab can carry is small, and saying so is more useful
than dressing it up. It is a \emph{single} run, reproducible from one
script but not repeated, with no fault injection: no peer was killed, no
partition induced, no message dropped. ``The three converge'' is
therefore an existence result and not a claim about convergence under
adversity, which repetitions and induced failures would be needed to
support. Its other boundaries are the ones §2 states --- an out-of-band
rendezvous, a fixed Director, trust-on-first-use certificates,
convergence completed by the framework's catch-up path. And by this
paper's own §8 the lab is measured against no baseline at all, since a
ported design has no distribution story to port. What it establishes is
that the domain's diff stays empty when the staging is genuinely
cross-machine, which is one more staging on the same axis as Lab A
rather than a second kind of evidence.

\textbf{Lab C --- the client.} With the stage held fixed, the client
varies: a person at a keyboard reading an on-screen grid; an automated
player sending moves over a pipe and reading a view it computes for
itself; two passive observers, one pulling the board and one receiving
it pushed; and a browser in which a player and a spectator are both
written in JavaScript and reach the game over the network. Six clients
over one \texttt{Well} --- some in C\#, one a PowerShell script, two in
a browser --- and the domain directory does not change for any of them.
Each adds an input adapter, an output adapter, or both. One of those
adapters was written by the client that reads through it rather than by
the author of the domain: the automated player is an instance of a large
language model and authored its own column-height view during the lab
(§3; disclosed in the acknowledgments).

Two counts belong here rather than one, because they answer different
questions and the larger flatters the smaller. Six is the number of
clients added, and it is the number that bears on the domain's diff:
each was a separate addition, and none of the six cost the domain an
edit. But six overstates how \emph{differently} the board was read.
Three of them --- the person at the keyboard, the observer that pulls,
the observer that receives a push --- read the same grid over three
different transports, and the browser's grid is that grid once more,
computed independently in JavaScript. Counted by representation instead
of by client, these labs hold \textbf{two}: a grid of cells, and a
vector of column heights. Counted by implementation, three: a C\#
renderer, a JavaScript one that mirrors it, and a PowerShell height
profile. The six is what supports the invariance claim; the two is the
honest measure of representational distance, and §3's argument that no
view is privileged rests on the two and not on the six. Both are
readable off one table, so a reader need not take either count on trust:

{\def\LTcaptype{none} % do not increment counter
\begin{longtable}[]{@{}
  >{\raggedright\arraybackslash}p{(\linewidth - 4\tabcolsep) * \real{0.3333}}
  >{\raggedright\arraybackslash}p{(\linewidth - 4\tabcolsep) * \real{0.3333}}
  >{\raggedright\arraybackslash}p{(\linewidth - 4\tabcolsep) * \real{0.3333}}@{}}
\toprule\noalign{}
\begin{minipage}[b]{\linewidth}\raggedright
Client
\end{minipage} & \begin{minipage}[b]{\linewidth}\raggedright
Reaches the domain by
\end{minipage} & \begin{minipage}[b]{\linewidth}\raggedright
Representation it reads
\end{minipage} \\
\midrule\noalign{}
\endhead
\bottomrule\noalign{}
\endlastfoot
console, driven by a person & keyboard, in process & grid of cells (C\#
renderer) \\
automated player & a pipe to a warm server & vector of column heights
(PowerShell) \\
observer & a query it issues, pulling & grid of cells (C\# renderer) \\
watch & a pushed frame & grid of cells (C\# renderer) \\
browser player & WebSocket & grid of cells (JavaScript) \\
browser spectator & server-sent events & grid of cells (JavaScript) \\
\end{longtable}
}

\textbf{Table 5.} The six clients of Lab C. Reading down the last column
gives the count that matters: two representations, one of them
implemented twice in different languages. Reading down the middle column
gives six distinct ways in, which is what the domain's empty diff was
measured against.

\textbf{Lab D --- the adapter.} One emitted frame reaches three
different projections, none of them the domain's. The board emits its
occupied cells --- a set it computes, in its own vocabulary
(\texttt{domain/Well.cs:322}, and §3 on why that is a fact it holds
rather than a view); a renderer lays that out as a terminal grid
(\texttt{actor/BoardRenderer.cs:19-48}); the browser produces the same
grid independently in JavaScript, its code commented as mirroring that
renderer (\texttt{web/Program.cs:169-178}); and the automated player
lifts the same frame into a column-height profile with its gaps, wells,
and aggregate figures (\texttt{tools/pile-scan.ps1:54-83}). Three
implementations over one fact, and no domain method for any of them ---
but two representations, not three, since the browser's grid mirrors the
terminal's; the height profile is the one that differs in kind.

\textbf{Lab E --- both sides testable, with the framework absent.} The
domain's own tests run with no staging bound at all: no host, no
transport, no sink, and no double --- not for what drives the domain and
not for what it emits. Stronger, and this is the form in which the claim
of §2 is checkable rather than interpretive: the framework is absent
from the domain's build graph altogether. The domain project is four
properties and no references ---

\textbf{The fence at the level of members.} Every verb-bearing type in
the domain is \texttt{internal}: the board, the pile, the shapes, the
piece types, the seven pieces, the rule exception
(\texttt{domain/Well.cs:34}, \texttt{domain/Pile.cs:21},
\texttt{domain/Frame.cs:29}, \texttt{domain/PieceType.cs:9},
\texttt{domain/Pieces.cs}, \texttt{domain/TetrisRuleException.cs:17}).
The library's whole public surface is one empty type, so not even the
actor's own project can invoke a domain operation in the host language
--- it obtains the assembly and passes it on. The compiler refuses where
reflection admits.

Where in-language access was genuinely wanted it had to be granted by
name, from inside the domain's own assembly, and it was: to the domain's
test project and to one console host
(\texttt{domain/AssemblyInfo.cs:8-9}). That is an authored exception,
enumerable, and extended to no other staging --- the four remaining
stagings of Experiment A cannot name a domain type at all. Two
qualifications belong with it. One grant is to a \emph{staging}, so the
claim is not that no staging may call the domain but that calling it
takes an act of authorship inside the domain, by name. And that grant is
no longer exercised --- the console drives the actor and names no domain
type, as its own project comment says --- so the fence is in practice
tighter than its source admits and the grant could be withdrawn without
breaking a build. It is reported rather than tidied away to purchase a
cleaner claim, and it is left standing: withdrawing it would tidy the
fence at the price of changing what this lab measured, which is a poor
trade for a count the paper leans on.

It would be wrong to say the domain references \emph{nothing}: it
targets a language runtime and uses that runtime's collections, as every
C\# project does. The precise claim, and the only one this paper uses,
is narrower --- \textbf{the domain declares no dependency on the
framework that runs it, or on any staging of it.}

\begin{Shaded}
\begin{Highlighting}[]
\NormalTok{\textless{}}\KeywordTok{Project} \OtherTok{Sdk=}\StringTok{"Microsoft.NET.Sdk"}\NormalTok{\textgreater{}}
\NormalTok{  \textless{}}\KeywordTok{PropertyGroup}\NormalTok{\textgreater{}}
\NormalTok{    \textless{}}\KeywordTok{TargetFramework}\NormalTok{\textgreater{}net9.0\textless{}/}\KeywordTok{TargetFramework}\NormalTok{\textgreater{}}
\NormalTok{    \textless{}}\KeywordTok{Nullable}\NormalTok{\textgreater{}enable\textless{}/}\KeywordTok{Nullable}\NormalTok{\textgreater{}}
\NormalTok{    \textless{}}\KeywordTok{RootNamespace}\NormalTok{\textgreater{}Tetris\textless{}/}\KeywordTok{RootNamespace}\NormalTok{\textgreater{}}
\NormalTok{    \textless{}}\KeywordTok{AssemblyName}\NormalTok{\textgreater{}TetrisDomain\textless{}/}\KeywordTok{AssemblyName}\NormalTok{\textgreater{}}
\NormalTok{  \textless{}/}\KeywordTok{PropertyGroup}\NormalTok{\textgreater{}}
\NormalTok{\textless{}/}\KeywordTok{Project}\NormalTok{\textgreater{}}
\end{Highlighting}
\end{Shaded}

--- with \texttt{dotnet\ list\ package} reporting no packages for it, no
framework namespace imported by any of its sources, no framework
attribute, and no framework interface implemented; the test project
references a test runner and the domain and nothing further. The suite
passes on that graph --- a regression check, as Lab A says, not a
coverage claim.

What this count does and does not separate is settled by the ported
baseline of Lab F, which has the same clean graph: zero project
references, zero packages, no framework named. So \emph{this} count
ties, and ``zero references, zero packages'' separates a clean domain
from an entangled one rather than one arrangement from the other. What
does not tie is on the driven side only: the ported application service
cannot be constructed without stand-ins for its three driven ports, and
twenty of its sixty-four tests cannot run without them, while the
journaled domain's tests need none.

{\def\LTcaptype{none} % do not increment counter
\begin{longtable}[]{@{}
  >{\raggedright\arraybackslash}p{(\linewidth - 4\tabcolsep) * \real{0.3333}}
  >{\raggedright\arraybackslash}p{(\linewidth - 4\tabcolsep) * \real{0.3333}}
  >{\raggedright\arraybackslash}p{(\linewidth - 4\tabcolsep) * \real{0.3333}}@{}}
\toprule\noalign{}
\begin{minipage}[b]{\linewidth}\raggedright
Lab
\end{minipage} & \begin{minipage}[b]{\linewidth}\raggedright
What varies
\end{minipage} & \begin{minipage}[b]{\linewidth}\raggedright
Measured
\end{minipage} \\
\midrule\noalign{}
\endhead
\bottomrule\noalign{}
\endlastfoot
A & 5 stagings (console, WebSocket, server-sent events, StageManager in
memory, StageManager over TLS) & 0 domain edits; the domain suite passes
unchanged (a regression check, not coverage) \\
B & 3 machines (containers, container-to-container TLS) & 0 domain edits
\textbf{and} 0 actor edits; three nodes converge, byte-identical
frame \\
C & 6 clients --- some in C\#, one a PowerShell script, two in a
browser; 2 distinct representations among them & 0 domain edits \\
D & 3 projections of one emitted frame & 0 domain methods for any
view \\
E & domain built and tested with the framework absent from its build
graph & 0 references, 0 packages, 0 doubles; the suite passes \\
G & the domain's own internal boundary --- one role cut into two,
stagings held still & 11 of 12 \textbf{hosts} unchanged (the 12th, one
line); 0 divergences over 47,783 steps; a whole game replayed into both
roles, game over included; record 2.32× (316 against 136), 6 extra
entries per landing \\
H & one reaction over 3 stagings' journals (single process, warm server
over a pipe, 3 containers read in place) & same 2 placements in the same
order in all three; the 3 containers' journals byte-identical, and the
acts matching verb for verb across stagings with entry ids offset by a
constant; 0 domain edits --- with no correlation handle, no landing act,
and entry ids not staging-invariant \\
I & the domain itself grows --- a score and a difficulty level added &
+30 lines of code (98 added lines in all, 62 doc comments and 6 blank),
0 code lines deleted (3 doc-comment lines removed), no signature or verb
changed; build graph byte-identical; 12/12 \textbf{hosts} keep running
at 0 edits; adoption costs 0 for four, 1 for five, 4 for each browser
host, 32 for the input host --- only where wanted \\
\end{longtable}
}

\textbf{Table 6.} What each lab measured. Every entry is a count taken
from the example's own history or build, and none of these rows is a
comparison --- the one comparison, against a ported arrangement built
for the purpose, is Lab F and is reported separately below.

\textbf{Lab F --- the ported baseline, built and counted.} The
comparison §9 rests on is a measurement rather than an estimate, and
this is where it was taken. An orthodox ports-and-adapters arrangement
of the same game was written --- its eleven rule files differ from the
journaled domain's by one line each, the namespace, so the comparison is
not rigged by writing a worse domain --- and four stagings were added to
it one commit at a time so that each could be counted: a console, a
WebSocket host, a REST-and-server-sent-events host, and an automated
player. All four run. Adding the baseline changed nothing on the
journaled side: the diff of the domain, the actor, the solution file,
and all fifteen existing hosts is empty, and the journaled side's tests
still pass.

A ported arrangement invites four estimates about what it costs a
domain, and measurement contradicts two of them, ties a third, and
leaves a difference that is not the one the estimates predict. Table 7
reports each against what was counted.

{\def\LTcaptype{none} % do not increment counter
\begin{longtable}[]{@{}
  >{\raggedright\arraybackslash}p{(\linewidth - 2\tabcolsep) * \real{0.5000}}
  >{\raggedright\arraybackslash}p{(\linewidth - 2\tabcolsep) * \real{0.5000}}@{}}
\toprule\noalign{}
\begin{minipage}[b]{\linewidth}\raggedright
Estimate
\end{minipage} & \begin{minipage}[b]{\linewidth}\raggedright
Measured against the built baseline
\end{minipage} \\
\midrule\noalign{}
\endhead
\bottomrule\noalign{}
\endlastfoot
an implementation per staging & \textbf{confirmed} --- 2, 1, 1, 2 across
the four stagings \\
a stand-in per driven port & \textbf{confirmed} --- three; the
application service cannot be constructed without them, and 20 of 64
tests cannot run at all \\
a domain edit per staging & \textbf{refuted} --- the WebSocket and the
server-sent-events stagings each cost the ported rule model nothing,
identical to the journaled side, twice over \\
a stand-in ``per side'' & \textbf{refuted} --- the driving side is zero,
because a hexagon \emph{implements} its driving port rather than
depending on it (§9) \\
a distinguishing build graph & \textbf{tie} --- both rule models: zero
project references, zero packages, no framework named \\
cost of a new \emph{capability} & cross-process persistence moved the
ported arrangement by 204 lines added and 9 removed, of which 61 changed
lines --- 56 added, 5 removed --- fall inside the rule model itself; the
journaled side gained the same capability, for the same client, at
nothing \\
\end{longtable}
}

\textbf{Table 7.} The ported estimates, measured. Two are refuted, one
ties, and the difference that survives is not the one the estimates
predict.

\textbf{What the baseline actually shows: the unit of account was
wrong.} A port costs a domain nothing per \emph{staging} and everything
per \emph{capability its ports do not already carry}. Two new transports
moved the hexagon not at all --- the same zero this paper reports,
reached by a different arrangement --- so invariance across transports
is not distinctive, and §2's counts should not be read as if it were.
What moved the hexagon was persistence: to let a game outlive its
process, the ported rule model had to grow a restore constructor and a
way to rebuild a pile from cells --- 61 changed lines, added and
removed, not added alone. The journaled arrangement acquired the same
capability, for the same client, at no cost to the domain, because a
record of the acts already existed. The measured difference is therefore
not \emph{per staging} but \emph{per capability of the kind a record
supplies for free}.

Read as a test of an arrangement rather than as a scoreboard, that is
the more useful thing the baseline does, and §9 leans on it there. This
lab holds two of the three constituents of the arrangement the paper
argues for --- a rule model free of dependencies, and a domain entered
by calls on its own operations rather than by calling outward --- and
lacks the third, declaring nothing about what becomes of what it emits.
Holding two of three, it produces neither surviving difference. Which is
what makes the paper's claim a claim about the assembly and not about
any part of it, since every part is credited to prior work.

One number in that table needs a caveat stated plainly, since it is the
most attackable figure in the baseline. Where a ported arrangement's
\emph{domain} ends is a matter of layout, and this baseline places its
rule model, its ports, and its application service in a single project
--- the ordinary arrangement at this size, and the one comparable to the
single project the journaled domain occupies. Read that way, persistence
cost 204 lines added and 9 removed. A reader who counts only the pure
rule model should read it as 56 added and 5 removed, and that is the
figure worth arguing with, being the part no choice of layout can move
outside the domain. Both are reported, either reading is available
without recomputing anything, and the choice changes no other
measurement in the table --- the stand-ins are required by whichever
project owns the port. On both readings the comparison is against
nothing.

\textbf{The rival design that would dissolve that number, and what it
costs instead.} There is a standard move in orthodox practice that would
make the 56 lines an artifact of one implementation choice rather than a
property of the arrangement: persist through a separate persistence
model and a mapper on the adapter side, so the rule model never learns
that anything is stored. If that were available here, the paper's
residual claim would fall with it. The reason it is not available lies
in the model rather than in the choice, and it can be read off the code.
\texttt{Pile} is \texttt{internal\ sealed} with a \texttt{private}
constructor (\texttt{domain/Pile.cs:21}, \texttt{:28}); the only
construction paths it offers are an empty pile and the landing
transition (\texttt{:35}, \texttt{:56}); and it enforces one invariant
by construction, that a pile never retains a complete row
(\texttt{:16}).

A mapper in another assembly therefore has three ways in, and each lands
somewhere. It can be handed a factory the domain authors, which is a
change to the rule model. It can be granted \texttt{InternalsVisibleTo}
--- which reaches \texttt{internal} members but not \texttt{private}
ones, so it still needs a domain-authored factory, and the grant is
itself a line inside the domain naming a persistence assembly. Or it can
use reflection, which reaches a private constructor and is what
mainstream object-relational mappers in fact do; that route asks nothing
of the domain, and it bypasses every check the constructors perform.

The baseline took the first road, and its own source records why. Its
rule model grew a factory,
\texttt{internal\ static\ Pile\ Of(int\ width,\ ImmutableHashSet\textless{}Position\textgreater{}\ cells)},
whose comment calls it the seam a state port needs to put a saved pile
back, added for the fourth staging because ``the only way in was through
the model'' (\texttt{baseline-hex/domain/model/Pile.cs:51}). It grew a
restore constructor whose comment names the cost outright --- the point
at which persistence reached the model (\texttt{Well.cs:75}). That
constructor re-asserts the invariants on the way in, the complete-row
check among them (\texttt{Well.cs:305}, \texttt{:314}). Its assembly
grants \texttt{InternalsVisibleTo} to its test suite and nothing else
(\texttt{baseline-hex/domain/AssemblyInfo.cs}), and both new members are
called from inside the domain project (\texttt{Well.cs:86},
\texttt{baseline-hex/domain/application/GameService.cs:197}), so the
mapper never touches the model. The separate-persistence-model design is
what this baseline does; it still required the seam.

Stated no more strongly than it should be: the claim is not that 56
lines is the necessary price. A model with public constructors and no
invariants would pay less, and would have less to protect. The claim is
that where a rule-bearing type is closed and enforces something, putting
a saved instance back is a domain-side act --- pay for it in the rule
model, as the baseline did, or reach past the constructors and forgo the
check.

The journaled side does neither, and the reason is qualitative rather
than a count. Reconstitution there is re-performing the recorded acts
through the domain's own operations, so \texttt{Integrate} runs again on
the way back in and the invariant holds on the reconstituted pile
exactly as it held on the original --- no seam, and nothing bypassed. A
mapper reconstitutes \emph{around} the invariant or is let in through a
door the domain opens; a replay reconstitutes \emph{through} it. What is
not measured is the third road: no reflection-based mapper variant was
built, so that comparison is an argument with checked premises rather
than a count, and building it is the next measurement this baseline
wants.

\textbf{One concession runs the other way, and belongs here rather than
in a footnote.} The build graph does not distinguish the arrangements: a
ported rule model is also free of project references, packages, and
framework names, so ``zero references, zero packages'' should stop
carrying comparative weight --- it distinguishes a clean domain from a
dirty one, not this arrangement from that one.

\textbf{A second concession was drafted beside it and did not survive
checking.} It read that the ported arrangement is cleaner on one point,
needing no counterpart to §2's empty public anchor type. But the
framework's seam takes an \emph{assembly}, not a type
(\texttt{Choreography/Theater/PerformanceV2.cs:60}), and reads one with
\texttt{GetTypes()} under a filter that admits internal types
(\texttt{Puppeteer/EventSourcing/DomainLibraries.cs:136,151}); the
framework's own CLI loads domain libraries by full path
(\texttt{PuppeteerCli/AttachCommand.cs:309}). Rebuilt with the anchor
deleted, the domain exposes zero public types, and a host holding no
reference to it, loading it by path, seeds a \texttt{Well} and spawns a
piece to the same reported state. The anchor is this example's
ergonomics and not the arrangement's cost, so the asymmetry claimed for
it is withdrawn.

\textbf{What remains unmeasured is distribution.} The two StageManager
stagings of Lab A and the three machines of Lab B have no ported
counterpart here, deliberately: a hexagon has no distribution story to
port, so building one would mean inventing an architecture rather than
measuring one. Lab B's zero therefore stands uncontested by measurement
rather than confirmed against a rival.

\textbf{Lab G --- re-decomposition.} Every other lab holds the domain
fixed and varies what surrounds it. This one does the opposite: it cuts
the domain's own internal boundary while the stagings stay where they
are. The \texttt{Well} was split into two roles, one owning the settled
pile and one owning the falling piece; the original's record was then
read --- by the same introspection read that rehydration performs, of
each action and its parameters --- and the new roles were driven to
perform their own. No journal was cut, transformed, or rewritten: the
original's acts are identical before and after, with nothing added or
removed, and both new records are append-only with contiguous entry
identifiers, each holding only its own role's verbs.

The exercise runs over a whole game rather than a prefix of one: 129
acts to game over, 29 landings, 136 entries, every emitted act present
in the record. Both roles reach the \texttt{Well}'s exact final state,
game over included --- which is the interesting part, because the re-cut
moved the authority for \emph{game over} to the pile role, and replaying
the record into the two reproduces it with both agreeing. Two things
make that a fair test rather than a private one. The framework's own
rehydration of the same journal reaches the same state, so the record is
a record in the ordinary sense and the lab's reading took nothing
special out of it. And the full length is available at the substrate
commit this paper cites, an earlier one having stopped the readable
record short of the game's end.

What was measured. Eleven of the twelve hosts were untouched while the
domain divided beneath them; the twelfth needed one line. The invariant
held over 60 games and 47,783 steps compared against the
single-\texttt{Well} version --- 9,265 landings, 3,336 rows cleared, 39
games ended, no divergence --- and 47,783 overlap checks that crossed
the new boundary, asking each role separately, found no overlap; the
roles never disagreed about a game being over. Those figures were re-run
on the corrected example at the substrate commit this paper cites, and
reproduced to the digit, so the equivalence result was never sensitive
to which state the example sat in; only the entry counts were. A random
player cleared no rows at all in twenty games, so two further playing
policies had to be written before the collapse was exercised. Nothing
became speculative, and two checks came out stronger than they had been.

What it cost. This is the one lab in which the domain does change, by
design: \texttt{Well.cs} was untouched, \texttt{Pile.cs} took two lines,
and 292 lines of new domain were written. Movement costs no cross-role
message; a landing costs two tells and two acknowledgements. The record
costs 2.32 times as much --- 316 entries across the two roles against
136 for the same game --- and almost all of the difference sits exactly
where the two roles have to speak. Of the 180 extra entries, 174 are the
29 landings at six entries each: a landing writes the piece's handover,
the pile's absorption, and a message and acknowledgement each way,
against the one act the undivided board wrote. One more is the second
seed, since two roles open two aggregates. The remaining five are
one-time declarations --- eleven templates across the two roles against
six for the one, each written once however many times its act recurs.
Entry counts depend on how verbs are recorded, and these were re-taken
after every command in the example became a parametrized Action: a
template plus a compact argument per act, where a literal script wrote
one full sentence per call. That encoding is why the original comes to
136 rather than the 130 an earlier run reported, and the ratio has moved
with it, from 2.38× through 2.29× to 2.32×. What the encoding does not
touch is the shape of the finding: 129 acts, 29 landings, six entries
per landing, and no divergence. A fourth state appeared, because landing
is now two acts by two parties rather than one; a guard had to be
restated; and a cell-set serialization moved inside the domain, a tell
carrying ordered scalars. One further consequence is larger than a cost,
and §8.4 states what it amounts to; the premises are here, each of them
checked. The undivided board built a complete frame by unioning the
pile's cells with the falling piece's and clipping to the interior, a
derivation performed inside the domain (\texttt{domain/Well.cs:322}). A
projection on the emitting plane is produced by a reaction belonging to
one actor, running on that reaction's thread once that actor's read lock
is released (\texttt{Puppeteer/IOutputSink.cs:106}), so it reaches that
actor's state and no other's. After the cut those two halves sit in two
actors, and neither role can push a whole frame. A third premise is a
fact about this lab's own artifact rather than about the substrate, and
is why nothing about the loss was measured: the split actor exposes no
sink, no formatter and no output wiring, so there was no push channel
here to lose in the first place. The reader that would join the two
records was not built and is not claimed.

The finding nobody anticipated. The record is not neutral between the
two new roles. Every act in the \texttt{Well}'s journal turned out to
belong to the piece role, and the pile role's own act --- absorbing
landed cells --- appears nowhere in the original, because under the
\texttt{Well} absorbing was never an act at all but a private
consequence of a tick. The pile role's record therefore does not
translate; it is generated, by the piece role performing and telling. So
a re-decomposition does not divide an existing record between two roles:
one inherits the original's voice, and the other's record is born as a
consequence of it. Whether that holds in general, a single case cannot
say.

\textbf{Lab H --- the same routine, recognized on three stagings.} One
reaction is defined outside the domain: it seeks a spawn, then the drop
that ends the piece's descent, and what it matches between them is a
placement. It is then run against three stagings --- the journal of a
single process; the journal a warm server keeps while a client drives it
over a named pipe; and the journals of the three containers, read inside
each node against its own store. It recognizes the same two placements,
in the same order, in all three. The three containers' journals are
byte-identical, and the domain's diff for the exercise is empty: a
reaction is defined outside the domain and the domain learns nothing of
it.

Three measured qualifications matter more than the confirmation. There
is \textbf{no correlation handle} in this record: a spawn names a piece
type and a drop names nothing, so the trajectory is tied by order alone.
There is \textbf{no landing act}: the board lands a piece privately,
inside a tick or a drop, so a piece that comes to rest under gravity
leaves its opening unclosed and the next piece's drop closes it --- two
placements matched at one closing entry, with the count right and the
correlation wrong, and nothing to signal it. And \textbf{entry
identifiers are not staging-invariant}: the three-node staging writes
three seeding acts where one process writes one, shifting everything
after, while the closing identifier is the only handle a match hands its
reader. The trajectory is short --- two placements --- so this is an
existence result about recognizability, not a measurement of it at
scale.

\textbf{Lab I --- the domain grows.} A score and a difficulty level were
added to the board: the first computable by any client and therefore
required to be single-valued, the second a rule over lines already
recorded. The growth was additive in the strict sense. Over the whole
domain directory the change is 98 lines added and 3 removed, across
three files rather than the two new ones alone: \texttt{Scoring.cs} at
39 lines and \texttt{Difficulty.cs} at 38, plus 21 added lines in
\texttt{Well.cs}. Of the 98 added, 30 are code; the other 68 are the
domain's own documentation convention --- 62 doc-comment lines and 6
blank --- and a reader should not take them for filler or for logic. All
3 removed lines are doc-comment prose in \texttt{Well.cs}. No code was
deleted, no signature changed and no verb removed, established by
diffing a sorted inventory of the board's members rather than by
inspection. The build graph came out byte-identical: no packages, no
references, and the two new files declare no imports whatsoever. The
domain's suite went from 44 tests to 60, the original 44 untouched.

The measurement this paper had never taken is the converse ripple. When
the domain grew, all twelve host projects kept running with no edit ---
verified by running them before and after from a throwaway worktree at
the pre-change commit, not merely by compiling them. The cost of
\emph{adopting} the new capability then fell only where it was wanted,
and the twelve divide exactly: nothing for four of them, one line for
five, four lines for each of the two browser hosts, and thirty-two for
the input host, the only place where the level changes what happens.
Those figures are the sum over both growth commits and both adoption
commits; taking only the first two understates the input host, which
acquires its thirty-two lines in the adoption commit rather than when
the level was added. That last is where the boundary this paper
predicted becomes visible: the domain owns what level the game is at,
the staging owns how fast gravity ticks, and the domain still holds no
clock --- which a reflection test now enforces.

One finding is a caution rather than a result, and it applies to every
replay measurement including this paper's. Three of the eight replay
fixtures did not \emph{fail}: they answered. A game recorded as having
cleared twenty rows replayed as having cleared two; one recorded at six
replayed at four; and the third, recorded at zero, replayed at zero ---
correct, because the acts it lost happened to clear nothing, so the
truncation was invisible in exactly the scalars a reader checks. The
cause was a defect in the substrate's index rather than anything about
the model, and it has since been fixed upstream, this lab being one of
two independent reports that found it. The lesson survives the fix, and
§8.4 states it as a rule: a replay checked only for internal consistency
will pass while being wrong, so this lab now carries the state recorded
at play time and asserts against it. That follows from what fusion means
and is close to a tautology; no fused version was built. Until that
exists, the comparison in this paragraph is an argument the paper makes,
not a result it reports.

The nine laboratories' source and datasets accompany this paper as
\texttt{paper09-data.zip}, laid out as the earlier papers in this series
lay theirs out: one directory per lab under \texttt{labs/paper09-labX-…}
holding its source and a README that states what it measures, what to
run and which figure of this paper it produces, and a directory of the
same name under \texttt{data/paper09-labX-…} holding its outputs and its
write-up. A suite-level \texttt{labs/paper09-README.md} answers what the
per-lab files cannot: that each lab stands alone and none depends on
another, which one to run if only one is run, the order by cost, and the
one interaction that can mislead --- Lab A measures that the domain did
\emph{not} change while Labs G and I change it deliberately, so running
Lab A's diff on their branches reports a difference that is a different
measurement rather than a failed one. Two of the nine have no source of
their own and say so in their README: Lab A, whose five stagings
\emph{are} hosts of the example and whose evidence is the empty diff
across them, and Lab H, whose reaction is quoted in its write-up. Lab F
is the largest, carrying the ported baseline in full. The labs are
count-based, so no large raw sample log is produced or omitted, and
every count cited above appears in full in the corresponding write-up.

Assembling that archive was itself a correction. The labs had been built
on separate branches of the examples repository and never merged, so the
material this section promises had been scattered across seven branches
and reachable only by knowing which; a reader could not have obtained it
from the repository named above. It is now consolidated, which is what
makes the promise in the previous paragraph one a reader can act on.

One correction made during this work belongs in plain sight, because it
touched the example, and an earlier draft of this appendix reported it
as done when the artifact this paper deposits did not yet carry it. The
frame each host pushes reaches its sink through a reaction over the
journal, and a reaction over the domain observes recorded
\emph{operations}, never literal script text --- a substrate rule, and a
deliberate one. The example's actor had been emitting its verbs as
literal text. So the push channel never engaged: no frame was written at
all, which cost three labs their evidence rather than merely changing
its form. The player hosts were unaffected, because each renders from
its own query of the board; the \emph{viewers} were the loss --- the
client-authored column-height view of §3 had no frame to read, the live
viewer of Lab D had nothing to repaint, and each node of Lab B wrote no
frame of its own.

The correction is now in the deposited artifact, and it is larger than
the one verb an earlier draft described. An operation is recorded as an
operation only when it carries an argument, so the seed carries the
board's dimensions and a spawn carries its piece; the five verbs that
move a falling piece carry nothing that varies --- there is no value in
\emph{move left} --- and each therefore carries the name of the act it
performs. The seam the actor speaks to no longer offers an
unparametrized form at all, which makes the omission unrepresentable
rather than merely discouraged, and the same conversion was applied to
the two-role actor of Lab G. Two consequences are reported where they
fall: entry counts move, because a template written once and a compact
argument per act is a different encoding from one full sentence per call
(§8.2), and Lab G's harness had been recovering the board's dimensions
by reading them out of the seed's text, so it now reads them from the
seed's arguments as it already did for a spawn.

The correction was to the example's actor, not to the domain: the
domain's diff across the entire period, including this fix, is empty ---
which is the invariance this appendix reports, holding through a change
to the machinery around it.

\subsection{Code provenance}\label{code-provenance}

Source-code references in this paper are paths in \textbf{this
repository}, and that is the one thing about provenance worth stating
plainly: the example the labs measure is vendored here, at
\texttt{labs/paper09-example/}, so the paper and the artifact it cites
are one git history and there is a single commit to mark rather than two
to reconcile. Example paths are cited relative to that directory ---
\texttt{domain/Well.cs:322}, \texttt{actor/IGameHost.cs:14} --- and the
ported baseline of Lab F, being a second and different arrangement,
keeps its own prefix: \texttt{baseline-hex/…} under
\texttt{labs/paper09-labF-ported-baseline/}. Line anchors resolve at the
commit this paper is deposited from; the example last moved at
\texttt{54a05ab}.

One thing a vendored copy cannot carry is a history, and §8 rests on
one: the dates that separate the domain's last commit from the stagings
added twenty-four days later. That chronology is checkable only where it
happened, in the repository the example was written in
(\texttt{github.com/alvaroNCubo/puppeteer-examples}), on \texttt{main}
at \texttt{ec9a2d4}: the domain's last commit is \texttt{fd8d94b}, the
three additions of 23 July are \texttt{33f151c}, \texttt{8d7536e} and
\texttt{974f62a}, and
\texttt{git\ diff\ fd8d94b..main\ -\/-\ Tetris/domain/} is empty across
the fifteen commits between. The copy here has one commit, the one that
placed it, which is why Lab A of Appendix A is the single lab a reader
cannot run against this repository and why its output is captured in
\texttt{data/} for a reader who would rather not clone anything.

Framework source is cited in five files and nowhere else, and only where
what is asserted is a property of the substrate that the example cannot
establish: the scope of the emitting plane, which §8's constraint on
admissible decompositions turns on (\texttt{Puppeteer/IOutputSink.cs},
\texttt{Puppeteer/StageHook.cs}); the attachment of a destination to a
live actor, which §9 uses to say when a destination is bound
(\texttt{Choreography/Theater/PerformanceV2.cs}); and what a domain
assembly must expose for the framework to read it, which Lab F's
withdrawn concession turns on and which an example carrying an anchor
cannot settle by itself
(\texttt{Puppeteer/EventSourcing/DomainLibraries.cs},
\texttt{PuppeteerCli/AttachCommand.cs}). Those five anchors resolve
against the \textbf{public} Puppeteer repository
(\texttt{github.com/alvaroNCubo/puppeteer}) at commit \texttt{2673d57},
which is also the minimum the labs' engine requirement names. An earlier
draft cited them against \texttt{dd67047}, a commit in the private Azure
repository the runtime is developed in --- a reference no reader could
resolve, and the substantive provenance defect this pass fixed. Every
other anchor in this paper is in the example, and every measurement
reported is taken on the example rather than on framework internals. The
example also contains two clients that no lab here counts --- a webcam
gesture input and an emulated constrained device --- left out because
their hardware legs were never run, so there is nothing measured to
report about them.

One thing about the labs' provenance belongs here rather than beside
each count, and it is more specific than a general disclaimer. This
draft's measurements were taken while the example was still moving, and
the labs do not all sit at one state of it: some were taken after the
actor correction reported at the end of Appendix A and some before it.
That correction changes how verbs are recorded, so it changes any figure
computed from entry counts. One such figure existed --- the record ratio
of Lab G --- and it has since been re-taken on the corrected example at
the substrate commit this paper cites, which moved it and left
everything else in that lab standing. No figure now reported comes from
an uncorrected state. The correction never touched the figures the
paper's argument rests on: a diff over a source directory, a count of
project references, a count of hosts that needed editing, and a
comparison of final board states are all indifferent to how an act was
written into the journal. Every lab is nonetheless to be re-verified at
the single published commit and reported at that commit as part of the
deposit, and where a count is sensitive to the example's state the
appendix says so rather than leaving a reader to reconcile dates.

A word on what is and is not public, since the two are easy to run
together. The repository named above is public; the two commits cited
above are not in it yet. They exist in the author's local clone and have
not been pushed, so a reader who visits that repository today will not
find them. Publishing them, pinning the resulting public commit here,
and archiving the snapshot under a Software Heritage persistent
identifier --- as the preceding paper in this series does --- belong to
the deposit of this version, and this section will carry the fixed
identifiers then. One cleanup was weighed for that same step and
deliberately not made: the domain's assembly carries an
\texttt{InternalsVisibleTo} grant to the console host which no host now
exercises (Lab E reports it). Withdrawing it would rest the fence on the
test suite's grant alone --- tidier, and a change to what Lab E measured
--- so it stays as reported. Until then the anchors are reproducible
from the accompanying archive, \texttt{paper09-data.zip}, whose contents
are listed above. The three-machine lab reproduces from a single script,
\texttt{Tetris/docker/run-demo.sh}, against a local container runtime.
The labs are count-based, so no large raw sample log is produced or
omitted.

\subsection{Acknowledgments}\label{acknowledgments}

The author used large language models (including Claude and ChatGPT) as
editorial assistants for language refinement, structural feedback, and
literature navigation. All original ideas, terminology, theoretical
constructs, and technical content presented in this work are solely the
author's.

One disclosure belongs with the labs rather than with the writing,
because a file it produced is cited above as evidence. The automated
player of Lab C is an instance of a large language model, driving the
game over a pipe; the view it reads the board through --- the
column-height projection of §3 (\texttt{tools/pile-scan.ps1}) --- was
written by that instance during the lab, with the author's permission,
to suit the form in which it reasons. That the adapter was authored by
the client rather than by the author of the domain is not incidental to
the argument. It is the sharpest case available of a projection
belonging to the party that reads it, and it is reported here so that a
reader can weigh the evidence knowing how it came about.

\subsection{References}\label{references}

Entries without a refereed venue --- practitioner articles, talks,
screencasts, and technical reports --- are given with their access date
and a permanent archival snapshot, and are marked as such where the
argument leans on them; where such a source sits behind a paywall, a
freely available presentation of the same idea is listed beside it.

The letter suffixes on this series' own entries denote the paper's
number in the series, not the order of appearance here ---
\texttt{2026c} is Paper 3, \texttt{2026d} Paper 4, \texttt{2026g} Paper
7, \texttt{2026h} Paper 8 --- so that a given paper carries the same
label wherever in the series it is cited. Letters absent below belong to
papers this one does not cite.

Agha, G. (1986). \emph{Actors: A model of concurrent computation in
distributed systems}. MIT Press.

Apel, S., Batory, D., Kästner, C., \& Saake, G. (2013).
\emph{Feature-oriented software product lines: Concepts and
implementation}. Springer.
\url{https://doi.org/10.1007/978-3-642-37521-7}

Armstrong, J. (2003). \emph{Making reliable distributed systems in the
presence of software errors} {[}Doctoral dissertation, Royal Institute
of Technology (KTH), Stockholm{]}.

Bernhardt, G. (2012a). \emph{Boundaries} {[}Conference talk, Software
Craftsmanship North America; practitioner source, no refereed venue,
freely available{]}. Destroy All Software.
\url{https://www.destroyallsoftware.com/talks/boundaries} (accessed 26
July 2026; archived at
\url{https://web.archive.org/web/20260703082615/https://www.destroyallsoftware.com/talks/boundaries})

Bernhardt, G. (2012b). \emph{Functional core, imperative shell}
{[}Screencast; practitioner source, no refereed venue, behind a paywall
--- see 2012a for the freely available presentation of the same idea{]}.
Destroy All Software.
\url{https://www.destroyallsoftware.com/screencasts/catalog/functional-core-imperative-shell}
(accessed 25 July 2026; archived at
\url{https://web.archive.org/web/20260723235945/https://www.destroyallsoftware.com/screencasts/catalog/functional-core-imperative-shell})

Brooks, F. P., Jr.~(1987). No silver bullet: Essence and accidents of
software engineering. \emph{Computer}, 20(4), 10--19.
\url{https://doi.org/10.1109/MC.1987.1663532}

Clements, P., \& Northrop, L. (2001). \emph{Software product lines:
Practices and patterns}. Addison-Wesley.

Cockburn, A. (2005). \emph{Hexagonal architecture} {[}Practitioner
article, no refereed venue{]}.
\url{https://alistair.cockburn.us/hexagonal-architecture/} (accessed 25
July 2026; archived at
\url{https://web.archive.org/web/20260708163140/https://alistair.cockburn.us/hexagonal-architecture})

Codd, E. F. (1970). A relational model of data for large shared data
banks. \emph{Communications of the ACM}, 13(6), 377--387.
\url{https://doi.org/10.1145/362384.362685}

Eugster, P. T., Felber, P. A., Guerraoui, R., \& Kermarrec, A.-M.
(2003). The many faces of publish/subscribe. \emph{ACM Computing
Surveys}, 35(2), 114--131. \url{https://doi.org/10.1145/857076.857078}

Evans, E. (2003). \emph{Domain-driven design: Tackling complexity in the
heart of software}. Addison-Wesley.

Fowler, M. (2004). \emph{Inversion of control containers and the
dependency injection pattern} {[}Practitioner article, no refereed
venue{]}. \url{https://martinfowler.com/articles/injection.html}
(accessed 25 July 2026; archived at
\url{https://web.archive.org/web/20260722213553/https://martinfowler.com/articles/injection.html})

Fowler, M. (2005). \emph{Event sourcing} {[}Practitioner article, no
refereed venue{]}.
\url{https://martinfowler.com/eaaDev/EventSourcing.html} (accessed 25
July 2026; archived at
\url{https://web.archive.org/web/20260714065500/https://martinfowler.com/eaaDev/EventSourcing.html})

Garlan, D., \& Notkin, D. (1991). Formalizing design spaces: Implicit
invocation mechanisms. In \emph{VDM '91: Formal Software Development
Methods}, Lecture Notes in Computer Science (Vol. 551, pp.~31--44).
Springer. \url{https://doi.org/10.1007/3-540-54834-3_5}

Gregor, S. (2006). The nature of theory in information systems.
\emph{MIS Quarterly}, 30(3), 611--642.

Helland, P. (2005). Data on the outside versus data on the inside.
\emph{Proceedings of the 2nd Biennial Conference on Innovative Data
Systems Research (CIDR '05)}, 144--153. (Reprinted in
\emph{Communications of the ACM}, 63(11), 111--118, 2020.
\url{https://doi.org/10.1145/3410623})

Helland, P. (2015). Immutability changes everything. \emph{Proceedings
of the 7th Biennial Conference on Innovative Data Systems Research (CIDR
'15)}. (Also published in \emph{ACM Queue}, 13(9), 2015,
\url{https://doi.org/10.1145/2857274.2884038}, and in
\emph{Communications of the ACM}, 59(1), 64--70, 2016,
\url{https://doi.org/10.1145/2844112})

Hewitt, C., Bishop, P., \& Steiger, R. (1973). A universal modular ACTOR
formalism for artificial intelligence. \emph{Proceedings of the 3rd
International Joint Conference on Artificial Intelligence (IJCAI)},
235--245.

Huang, P., Guo, C., Zhou, L., Lorch, J. R., Dang, Y., Chintalapati, M.,
\& Yao, R. (2017). Gray failure: The Achilles' heel of cloud-scale
systems. \emph{Proceedings of the 16th Workshop on Hot Topics in
Operating Systems (HotOS '17)}, Whistler, BC, Canada.
\url{https://doi.org/10.1145/3102980.3103005}

Johnson, R. E., \& Foote, B. (1988). Designing reusable classes.
\emph{Journal of Object-Oriented Programming}, 1(2), 22--35.

Kaldor, J., Mace, J., Bejda, M., Gao, E., Kuropatwa, W., O'Neill, J.,
Ong, K. W., Schaller, B., Shan, P., Viscomi, B., Venkataraman, V.,
Veeraraghavan, K., \& Song, Y. J. (2017). Canopy: An end-to-end
performance tracing and analysis system. \emph{Proceedings of the 26th
Symposium on Operating Systems Principles (SOSP '17)}, 34--50.
\url{https://doi.org/10.1145/3132747.3132749}

Kleppe, A., Warmer, J., \& Bast, W. (2003). \emph{MDA explained: The
model driven architecture --- Practice and promise}. Addison-Wesley.

Kleppmann, M., \& Kreps, J. (2015). Kafka, Samza and the Unix philosophy
of distributed data. \emph{IEEE Data Engineering Bulletin}, 38(4),
4--14.

Kratzke, N., \& Quint, P.-C. (2017). Understanding cloud-native
applications after 10 years of cloud computing --- A systematic mapping
study. \emph{Journal of Systems and Software}, 126, 1--16.

Lamport, L. (1978). Time, clocks, and the ordering of events in a
distributed system. \emph{Communications of the ACM}, 21(7), 558--565.
\url{https://doi.org/10.1145/359545.359563}

Laurel, B. (1991). \emph{Computers as theatre}. Addison-Wesley. (2nd
ed., Addison-Wesley Professional, 2013.)

Mace, J., Roelke, R., \& Fonseca, R. (2015). Pivot tracing: Dynamic
causal monitoring for distributed systems. \emph{Proceedings of the 25th
Symposium on Operating Systems Principles (SOSP '15)}, 378--393.
\url{https://doi.org/10.1145/2815400.2815415}

Martin, R. C. (2017). \emph{Clean architecture: A craftsman's guide to
software structure and design}. Prentice Hall.

Mellor, S. J., \& Balcer, M. J. (2002). \emph{Executable UML: A
foundation for model-driven architecture}. Addison-Wesley.

Meyer, B. (1988). \emph{Object-oriented software construction}. Prentice
Hall.

Moseley, B., \& Marks, P. (2006). \emph{Out of the tar pit} {[}Presented
at Software Practice Advancement (SPA 2006); no formal proceedings{]}.
\url{https://curtclifton.net/papers/MoseleyMarks06a.pdf} (accessed 25
July 2026; archived at
\url{https://web.archive.org/web/20260718134750/https://curtclifton.net/papers/MoseleyMarks06a.pdf})

Murphy, G. C., Notkin, D., \& Sullivan, K. (1995). Software reflexion
models: Bridging the gap between source and high-level models.
\emph{Proceedings of the 3rd ACM SIGSOFT Symposium on the Foundations of
Software Engineering (SIGSOFT '95)}; \emph{ACM SIGSOFT Software
Engineering Notes}, 20(4), 18--28.
\url{https://doi.org/10.1145/222132.222136}

Murphy, G. C., Notkin, D., \& Sullivan, K. J. (2001). Software reflexion
models: Bridging the gap between design and implementation. \emph{IEEE
Transactions on Software Engineering}, 27(4), 364--380. (The extended
journal version of the 1995 paper, and the canonical citable form.)

Overeem, M., Spoor, M., \& Jansen, S. (2017). The dark side of event
sourcing: Managing data conversion. \emph{Proceedings of the 24th IEEE
International Conference on Software Analysis, Evolution and
Reengineering (SANER '17)}, 193--204.
\url{https://doi.org/10.1109/SANER.2017.7884621}

Overeem, M., Spoor, M., Jansen, S., \& Brinkkemper, S. (2021). An
empirical characterization of event sourced systems and their schema
evolution --- Lessons from industry. \emph{Journal of Systems and
Software}, 178, 110970.

Parnas, D. L. (1972). On the criteria to be used in decomposing systems
into modules. \emph{Communications of the ACM}, 15(12), 1053--1058.
\url{https://doi.org/10.1145/361598.361623}

Passos, L., Terra, R., Valente, M. T., Diniz, R., \& Mendonça, N.
(2010). Static architecture-conformance checking: An illustrative
overview. \emph{IEEE Software}, 27(5), 82--89.

Perry, D. E., \& Wolf, A. L. (1992). Foundations for the study of
software architecture. \emph{ACM SIGSOFT Software Engineering Notes},
17(4), 40--52. \url{https://doi.org/10.1145/141874.141884}

Peyton Jones, S., \& Wadler, P. (1993). Imperative functional
programming. \emph{Proceedings of the 20th ACM SIGPLAN-SIGACT Symposium
on Principles of Programming Languages (POPL '93)}, 71--84.
\url{https://doi.org/10.1145/158511.158524}

Rivera, A. (2026c). Reactions and the partition: opt-in eventual
consistency in actor-native systems. \emph{Puppeteer Papers Series},
Paper 3 {[}Preprint{]}. Zenodo.
\url{https://doi.org/10.5281/zenodo.20792156}

Rivera, A. (2026d). Preserving semantic continuity across actors: a
tell-based approach without orchestration. \emph{Puppeteer Papers
Series}, Paper 4 {[}Preprint{]}. Zenodo.
\url{https://doi.org/10.5281/zenodo.21207062}

Rivera, A. (2026g). After the substrate: building software without a
datacenter. \emph{Puppeteer Papers Series}, Paper 7 {[}Preprint{]}.
Zenodo. \url{https://doi.org/10.5281/zenodo.20398998}

Rivera, A. (2026h). Inference without Authority: the three authorities
that govern an output. \emph{Puppeteer Papers Series}, Paper 8
{[}Preprint{]}. Zenodo. \url{https://doi.org/10.5281/zenodo.21499637}

Sambasivan, R. R., Shafer, I., Mace, J., Sigelman, B. H., Fonseca, R.,
\& Ganger, G. R. (2016). Principled workflow-centric tracing of
distributed systems. \emph{Proceedings of the 7th ACM Symposium on Cloud
Computing (SoCC '16)}, 401--414.
\url{https://doi.org/10.1145/2987550.2987568}

Shaw, M. (2003). Writing good software engineering research papers.
\emph{Proceedings of the 25th International Conference on Software
Engineering (ICSE '03)}, 726--736.

Sigelman, B. H., Barroso, L. A., Burrows, M., Stephenson, P., Plakal,
M., Beaver, D., Jaspan, S., \& Shanbhag, C. (2010). \emph{Dapper, a
large-scale distributed systems tracing infrastructure} (Technical
Report dapper-2010-1) {[}Technical report, no refereed venue{]}. Google
Inc.~\url{https://research.google/pubs/dapper-a-large-scale-distributed-systems-tracing-infrastructure/}
(accessed 25 July 2026; archived at
\url{https://web.archive.org/web/20260705194849/https://research.google/pubs/dapper-a-large-scale-distributed-systems-tracing-infrastructure/})

Stol, K.-J., \& Fitzgerald, B. (2018). The ABC of software engineering
research. \emph{ACM Transactions on Software Engineering and
Methodology}, 27(3), Article 11. \url{https://doi.org/10.1145/3241743}

Sullivan, K. J., \& Notkin, D. (1992). Reconciling environment
integration and software evolution. \emph{ACM Transactions on Software
Engineering and Methodology}, 1(3), 229--268.

Terra, R., \& Valente, M. T. (2009). A dependency constraint language to
manage object-oriented software architectures. \emph{Software: Practice
and Experience}, 39(12), 1073--1094.
\url{https://doi.org/10.1002/spe.931}

Terry, D. B., Demers, A. J., Petersen, K., Spreitzer, M. J., Theimer, M.
M., \& Welch, B. B. (1994). Session guarantees for weakly consistent
replicated data. \emph{Proceedings of the Third International Conference
on Parallel and Distributed Information Systems (PDIS '94)}, 140--149.

Waldo, J., Wyant, G., Wollrath, A., \& Kendall, S. (1994). \emph{A note
on distributed computing} (Technical Report SMLI TR-94-29). Sun
Microsystems Laboratories. (Reprinted in J. Vitek \& C. Tschudin, Eds.,
\emph{Mobile object systems: Towards the programmable internet}, Lecture
Notes in Computer Science, Vol. 1222, pp.~49--64. Springer, 1997.
\url{https://doi.org/10.1007/3-540-62852-5_6})

Wiggins, A. (2011). \emph{The twelve-factor app} {[}Practitioner source,
no refereed venue{]}. \url{https://12factor.net} (accessed 25 July 2026;
archived at
\url{https://web.archive.org/web/20260724222316/https://12factor.net/})

Young, G. (2010). \emph{CQRS documents} {[}Practitioner source, no
refereed venue{]}.
\url{https://cqrs.files.wordpress.com/2010/11/cqrs_documents.pdf}
(accessed 25 July 2026; archived at
\url{https://web.archive.org/web/20240317190715/https://cqrs.files.wordpress.com/2010/11/cqrs_documents.pdf})

\end{document}
